\documentclass[11pt]{article}
\input{../shared/project_style.tex}
\rhead{Module 1 Physics Note}

\begin{document}
\DocTitle{Module 1: Stellarator Equilibrium and Magnetic Geometry}{Textbook Physics Note, Version 1.0}{From elementary magnetic confinement to three-dimensional flux coordinates and equilibrium data products}

\begin{overviewbox}
\textbf{Purpose and status.} This chapter develops the physics and mathematics needed to understand a stellarator equilibrium from the beginning. It assumes no prior plasma-physics background beyond elementary calculus and vectors. Every abbreviation is expanded at first use, every coordinate and field variable is defined, and the physical meaning of the principal equations is stated alongside the mathematics. The chapter explains what an equilibrium is, how magnetic surfaces and field lines are represented, what rotational transform and magnetic shear mean, why a stellarator requires a genuinely three-dimensional description, and which geometric quantities must be delivered to the remaining kinetic-magnetohydrodynamic modules. This is a completed educational physics note. It does \emph{not} claim that a validated three-dimensional equilibrium solver has been implemented in this repository; the corresponding implementation note remains a solver plan.
\end{overviewbox}

\ProjectTOC

\section*{How to use this chapter}
\addcontentsline{toc}{section}{How to use this chapter}
A reader encountering plasma confinement for the first time should read Sections~\ref{sec:physical_picture}--\ref{sec:equilibrium_equations} in order. Sections~\ref{sec:flux_surfaces}--\ref{sec:magnetic_coordinates} introduce the geometry that is required by Module~4, which computes Alfv\'en continua and eigenmodes. Sections~\ref{sec:fourier_geometry}--\ref{sec:equilibrium_solvers} explain how three-dimensional stellarator equilibria are represented and computed. Sections~\ref{sec:outputs}--\ref{sec:verification} define the data products and verification obligations for the project.

The central dependency chain is
\begin{equation}
\begin{aligned}
\text{coils and plasma profiles}
&\longrightarrow \text{magnetohydrodynamic force balance}\\
&\longrightarrow \text{magnetic surfaces and field-line geometry}\\
&\longrightarrow \text{orbit, wave, and transport calculations}.
\end{aligned}
\end{equation}
The purpose of Module~1 is to establish the first three links reliably enough that every downstream calculation uses a common geometry and a common set of conventions.

\begin{physicsbox}
\textbf{Learning objectives.} After completing this chapter, the reader should be able to:
\begin{enumerate}[leftmargin=1.6em]
\item explain why toroidal magnetic confinement requires helical field-line winding;
\item derive the static ideal-magnetohydrodynamic force-balance equation and interpret magnetic pressure and tension;
\item define magnetic surfaces, toroidal and poloidal flux, rotational transform, safety factor, and magnetic shear;
\item distinguish geometric angles from straight-field-line magnetic coordinates;
\item interpret three-dimensional stellarator Fourier coefficients and field periods;
\item describe fixed-boundary and free-boundary equilibrium calculations;
\item identify the geometry, derivatives, metrics, and metadata required by the downstream orbit, wave, transport, and surrogate modules; and
\item design verification tests that separate numerical convergence from physical validation.
\end{enumerate}
\end{physicsbox}

\section*{Recommended prerequisite basics notes}
\addcontentsline{toc}{section}{Recommended prerequisite basics notes}
This chapter is designed to be readable without first completing the entire basics curriculum. The following notes provide useful parallel explanations.

\begin{longtable}{>{\RaggedRight\arraybackslash}p{0.15\textwidth} >{\RaggedRight\arraybackslash}p{0.28\textwidth} >{\RaggedRight\arraybackslash}p{0.47\textwidth}}
\toprule
\textbf{Basics note} & \textbf{Subject} & \textbf{Connection to Module 1}\\
\midrule
\endfirsthead
\toprule
\textbf{Basics note} & \textbf{Subject} & \textbf{Connection to Module 1}\\
\midrule
\endhead
\href{run:basics/basics_01_units_dimensions_and_normalization.pdf}{Basics 01} & Units, dimensions, and normalization & Reviews the International System of Units, dimensional checks, characteristic scales, and normalized variables.\\
\href{run:basics/basics_02_vectors_fields_and_coordinate_systems.pdf}{Basics 02} & Vectors, fields, and coordinates & Introduces vector components, basis vectors, cylindrical coordinates, and curvilinear coordinates.\\
\href{run:basics/basics_03_vector_calculus_for_plasma_physics.pdf}{Basics 03} & Vector calculus & Reviews gradient, divergence, curl, flux, and integral theorems.\\
\href{run:basics/basics_11_electric_and_magnetic_fields.pdf}{Basics 11} & Electric and magnetic fields & Introduces electromagnetic fields and their physical interpretation.\\
\href{run:basics/basics_12_maxwell_equations.pdf}{Basics 12} & Maxwell equations & Gives the field equations from which the magnetostatic relations used here are obtained.\\
\href{run:basics/basics_17_what_is_a_plasma.pdf}{Basics 17} & What is a plasma? & Explains quasineutrality, collective behavior, species, and plasma parameters.\\
\href{run:basics/basics_21_single_fluid_and_multifluid_descriptions.pdf}{Basics 21} & Fluid descriptions & Explains how particle species are represented by density, velocity, and pressure fields.\\
\href{run:basics/basics_24_derivation_of_ideal_magnetohydrodynamics.pdf}{Basics 24} & Ideal magnetohydrodynamics & Derives the bulk equations whose time-independent limit defines the present equilibrium problem.\\
\href{run:basics/basics_26_ideal_ohms_law_and_flux_freezing.pdf}{Basics 26} & Ideal Ohm's law and flux freezing & Explains why magnetic field lines move with a highly conducting ideal plasma.\\
\href{run:basics/basics_59_magnetic_confinement_concepts.pdf}{Basics 59} & Magnetic-confinement concepts & Provides a shorter conceptual introduction to toroidal confinement.\\
\href{run:basics/basics_60_tokamaks_and_stellarators.pdf}{Basics 60} & Tokamaks and stellarators & Compares the two principal toroidal magnetic-confinement concepts.\\
\href{run:basics/basics_61_toroidal_coordinates_and_magnetic_surfaces.pdf}{Basics 61} & Toroidal coordinates and magnetic surfaces & Develops the geometric language used in Sections~\ref{sec:flux_surfaces}--\ref{sec:magnetic_coordinates}.\\
\href{run:basics/basics_62_rotational_transform_magnetic_shear_and_resonances.pdf}{Basics 62} & Rotational transform and shear & Expands the field-line winding concepts used throughout the project.\\
\href{run:basics/basics_63_stellarator_symmetry_and_fourier_geometry.pdf}{Basics 63} & Stellarator symmetry and Fourier geometry & Gives a focused account of the spectral surface representation used in Section~\ref{sec:fourier_geometry}.\\
\href{run:basics/basics_65_magnetic_field_coils_and_biot_savart_law.pdf}{Basics 65} & Coils and the Biot--Savart law & Connects engineering coil currents to the vacuum magnetic field.\\
\href{run:basics/basics_66_mhd_equilibrium_codes_and_vmec_concepts.pdf}{Basics 66} & Equilibrium codes and VMEC concepts & Gives a companion introduction to numerical three-dimensional equilibrium solvers.\\
\bottomrule
\caption{Recommended supporting notes. Basics 59--66 are curriculum scaffolds in the present repository baseline; Module 1 supplies the complete first-pass explanation needed here.}
\label{tab:module1_prerequisites}
\end{longtable}

\SymbolsAcronymsSection
\begin{symtable}
AE & Alfv\'en eigenmode & A coherent oscillation whose restoring force is largely magnetic tension; computed in Module 4.\\
Boozer coordinates & Straight-field-line magnetic coordinates & Flux coordinates in which magnetic-field representations are especially useful for particle orbits and neoclassical transport.\\
Cylindrical coordinates & $(R,\phi,Z)$ & $R$ is distance from the symmetry axis, $\phi$ is the geometric toroidal angle, and $Z$ is vertical position.\\
EP & Energetic particle & A particle population with energy well above the bulk thermal energy, such as fusion-born alpha particles or neutral-beam ions.\\
FEM & Finite-element method & A numerical method based on weak forms and piecewise basis functions.\\
FFT & Fast Fourier transform & An efficient algorithm for converting sampled periodic data to Fourier coefficients and back.\\
LCFS & Last closed flux surface & The outermost closed magnetic surface used to define the confined plasma region in a nested-surface model.\\
MHD & Magnetohydrodynamics & A fluid description in which a conducting plasma interacts with electric and magnetic fields.\\
NBI & Neutral beam injection & A heating and current-drive method that injects fast neutral atoms, which ionize inside the plasma.\\
PINN & Physics-informed neural network & A neural approximation constrained partly by differential-equation residuals and boundary conditions.\\
QA & Quasi-axisymmetry & A form of quasisymmetry in which magnetic-field strength is approximately axisymmetric in suitable magnetic coordinates.\\
QH & Quasi-helical symmetry & A form of quasisymmetry in which magnetic-field strength approximately follows a helical symmetry.\\
SI & International System of Units & The unit system used unless a normalized variable is explicitly introduced.\\
VMEC & Variational Moments Equilibrium Code & A widely used code that computes ideal-MHD equilibria with nested magnetic surfaces using a variational formulation.\\
$\mathbf x$ & Position vector & A point in physical three-dimensional space, usually measured in metres.\\
$(x,y,z)$ & Cartesian coordinates & Orthogonal coordinates used for general vector operations.\\
$(R,\phi,Z)$ & Cylindrical coordinates & Coordinates adapted to toroidal devices; $\phi$ increases around the major axis.\\
$(s,\theta,\zeta)$ & Magnetic flux coordinates & $s$ labels a magnetic surface, $\theta$ is a poloidal angle, and $\zeta$ is a toroidal angle.\\
$R_0$ & Reference major radius & A characteristic distance from the device symmetry axis to the magnetic axis.\\
$a$ & Reference minor radius & A characteristic cross-sectional radius of the plasma.\\
$A$ & Aspect ratio & The dimensionless ratio $A=R_0/a$.\\
$N_{\mathrm{fp}}$ & Number of field periods & Number of repeated toroidal sectors in a stellarator magnetic configuration.\\
$\mathbf B$ & Magnetic flux density & The magnetic field measured in tesla; $B=|\mathbf B|$ is its magnitude.\\
$\mathbf b$ & Magnetic-field unit vector & $\mathbf b=\mathbf B/B$.\\
$\mathbf E$ & Electric field & Force per unit charge, measured in volts per metre.\\
$\mathbf J$ & Electric current density & Electric current per unit area, measured in amperes per square metre.\\
$p$ & Scalar plasma pressure & Isotropic thermodynamic pressure, measured in pascals, in the ideal-MHD equilibrium model.\\
$\rho_m$ & Mass density & Plasma mass per unit volume, measured in kilograms per cubic metre.\\
$n_s$ & Number density of species $s$ & Number of particles of species $s$ per unit volume.\\
$T_s$ & Temperature of species $s$ & Thermal energy scale; when stated in kelvin it is multiplied by Boltzmann's constant in energy expressions.\\
$q_s$ & Charge of species $s$ & Electric charge of a particle species.\\
$m_s$ & Mass of species $s$ & Particle mass of species $s$.\\
$\mu_0$ & Vacuum permeability & Magnetic constant, approximately $4\pi\times10^{-7}\ \mathrm{H\,m^{-1}}$.\\
$\epsilon_0$ & Vacuum permittivity & Electric constant in Maxwell's equations.\\
$\psi$ & Toroidal-flux coordinate & A magnetic-surface label proportional to toroidal magnetic flux; the exact normalization must always be stated.\\
$\Psi_t$ & Toroidal magnetic flux & Magnetic flux through a poloidal cross-section bounded by a magnetic surface.\\
$\Psi_p$ & Poloidal magnetic flux & Magnetic flux through a toroidal ribbon bounded by a magnetic surface.\\
$s$ & Normalized flux coordinate & Commonly $s=\Psi_t/\Psi_{t,a}$, where the subscript $a$ denotes the plasma boundary.\\
$\rho$ & Normalized radius & Often $\rho=\sqrt{s}$; this symbol is not the mass density $\rho_m$.\\
$\theta$ & Poloidal magnetic angle & Periodic angle that advances around a short way through the plasma cross-section.\\
$\zeta$ & Toroidal magnetic angle & Periodic angle that advances around the long way through the torus.\\
$\iota$ & Rotational transform & Average poloidal turns per toroidal turn made by a magnetic field line.\\
$q$ & Safety factor & In the convention used here, $q=1/\iota$; common in tokamak literature.\\
$\hat s$ & Magnetic shear & A dimensionless measure of radial variation of field-line pitch; a specific definition is given in Section~\ref{sec:iota}.\\
$m$ & Poloidal mode number & Integer counting variation in the poloidal angle.\\
$n$ & Toroidal mode number & Integer counting variation in the toroidal angle under the stated Fourier convention.\\
$\sqrt{g}$ & Coordinate Jacobian & Volume scale factor satisfying $dV=\sqrt{g}\,ds\,d\theta\,d\zeta$.\\
$g_{ij}$ & Covariant metric tensor & Dot products of covariant basis vectors: $g_{ij}=\partial_i\mathbf x\cdot\partial_j\mathbf x$.\\
$g^{ij}$ & Contravariant metric tensor & Dot products of gradients: $g^{ij}=\nabla u^i\cdot\nabla u^j$, where $u^i$ denotes a coordinate.\\
$V(s)$ & Enclosed volume & Physical volume inside the surface labeled by $s$.\\
$V'(s)$ & Differential volume & Derivative $dV/ds$, used in flux-surface transport equations.\\
$\langle A\rangle_s$ & Flux-surface average & Surface-weighted average of a scalar $A$ over the magnetic surface $s$.\\
$\boldsymbol\kappa$ & Magnetic curvature & $\boldsymbol\kappa=\mathbf b\cdot\nabla\mathbf b$, directed toward the local center of curvature of a field line.\\
$\beta$ & Plasma beta & Ratio of plasma pressure to magnetic pressure; a precise local or volume-averaged definition must be stated.\\
$W$ & MHD energy functional & Total magnetic plus internal energy used in a variational equilibrium formulation.\\
$\gamma_{\mathrm{ad}}$ & Adiabatic index & Ratio of specific heats in an ideal-gas pressure model; not an instability growth rate.\\
$\omega$ & Angular frequency & Oscillation rate in radians per second.\\
$k_\parallel$ & Parallel wave number & Spatial phase variation per unit length along $\mathbf B$.\\
\end{symtable}


% The chapter is modularized so each major block can be audited and compiled independently.
% Module 1 section file. Included by ../module1_stellarator_equilibrium_geometry.tex.
\section{The physical problem: confining a hot plasma}
\label{sec:physical_picture}

\subsection{What is being confined?}
A plasma is an ionized gas containing positively charged ions, negatively charged electrons, and sometimes neutral particles. In a fusion plasma the temperature is so high that ordinary material walls cannot directly contain the central plasma without being rapidly cooled or damaged. Magnetic confinement instead uses the Lorentz force
\begin{equation}
\mathbf F_s=q_s\left(\mathbf E+\mathbf v_s\times\mathbf B\right)
\label{eq:lorentz_force}
\end{equation}
to influence charged-particle motion. Here $\mathbf F_s$ is the force on a particle of species $s$, $q_s$ is its charge, $\mathbf v_s$ is its velocity, $\mathbf E$ is the electric field, and $\mathbf B$ is the magnetic flux density.

The magnetic part of the Lorentz force is perpendicular to the instantaneous particle velocity. It therefore bends a charged particle's path without directly changing its kinetic energy. In a sufficiently strong and slowly varying magnetic field, a particle executes rapid gyromotion around a magnetic field line while moving more freely along the field. This anisotropy is the basic reason magnetic fields can reduce transport across the field.

\begin{physicsbox}
\textbf{Central confinement idea.} A strong magnetic field does not place particles at fixed points. It organizes their trajectories so that cross-field motion is much slower than motion along the field. Successful confinement therefore requires field lines themselves to remain inside the device for long distances and, ideally, to lie on closed toroidal surfaces.
\end{physicsbox}

\subsection{Why close the magnetic field into a torus?}
If a magnetic field line ends on a material wall, particles streaming along it can be lost rapidly. A straight magnetic bottle can reflect some particles using stronger magnetic fields at its ends, but end losses remain a major challenge. Toroidal devices bend the plasma and magnetic field into a ring so that field lines need not terminate at physical end plates.

A torus has two natural directions:
\begin{itemize}[leftmargin=1.6em]
\item the \textbf{toroidal direction}, which goes the long way around the ring;
\item the \textbf{poloidal direction}, which goes the short way around a cross-section of the ring.
\end{itemize}
A useful confining field contains both toroidal and poloidal components. Their combination makes a field line wind helically around the torus instead of repeatedly sampling only one side of the device.

\subsection{Tokamak and stellarator strategies}
A \textbf{tokamak} produces most of its toroidal field using external coils and a substantial part of its poloidal field using a large toroidal plasma current. A \textbf{stellarator} produces the required three-dimensional field-line twist primarily through external coils and shaped magnetic geometry. The stellarator therefore does not require a large transformer-driven plasma current for its fundamental confinement geometry.

This difference has several consequences.
\begin{enumerate}[leftmargin=1.6em]
\item A stellarator equilibrium is intrinsically three-dimensional: magnetic quantities generally depend on both poloidal and toroidal angles.
\item The rotational transform can be sustained in steady state by external coils.
\item Coil and boundary shape become direct design variables for confinement, stability, and energetic-particle behavior.
\item Numerical representation is more demanding because axisymmetry cannot usually reduce the problem to two dimensions.
\end{enumerate}

\begin{notebox}
\textbf{Equilibrium is not confinement quality.} A magnetic configuration may satisfy force balance yet have poor particle confinement, undesirable neoclassical transport, large fast-ion losses, or unstable waves. Module~1 determines the equilibrium geometry. Modules~3--7 determine important consequences of that geometry.
\end{notebox}

\section{From particles to magnetohydrodynamics}
\label{sec:mhd_model}

\subsection{Why use a fluid model?}
A kinetic description follows distribution functions in six-dimensional phase space: three position coordinates and three velocity coordinates. That description is essential for resonant energetic particles and is developed later in the project. The bulk thermal plasma, however, can often be represented on large spatial and temporal scales using fluid variables such as mass density $\rho_m(\mathbf x,t)$, flow velocity $\mathbf u(\mathbf x,t)$, and pressure $p(\mathbf x,t)$.

Magnetohydrodynamics, abbreviated MHD, combines fluid dynamics with Maxwell's equations. The ideal-MHD model assumes, among other approximations, high electrical conductivity, a single bulk fluid velocity, and usually an isotropic scalar pressure. These assumptions are not universally valid, but they provide the standard starting point for macroscopic equilibrium and stability; a systematic development is given by Freidberg~\cite{freidberg2014}.

\subsection{Ideal-MHD equations before taking the equilibrium limit}
A common ideal-MHD system is
\begin{align}
\frac{\partial \rho_m}{\partial t}+\nabla\cdot(\rho_m\mathbf u)&=0,
\label{eq:mass_conservation}\\
\rho_m\left(\frac{\partial \mathbf u}{\partial t}+\mathbf u\cdot\nabla\mathbf u\right)&=-\nabla p+\mathbf J\times\mathbf B,
\label{eq:mhd_momentum}\\
\frac{\partial \mathbf B}{\partial t}&=\nabla\times(\mathbf u\times\mathbf B),
\label{eq:induction}\\
\nabla\cdot\mathbf B&=0,
\label{eq:divB}\\
\nabla\times\mathbf B&=\mu_0\mathbf J,
\label{eq:ampere_mhd}\\
\frac{d}{dt}\left(\frac{p}{\rho_m^{\gamma_{\mathrm{ad}}}}\right)&=0.
\label{eq:adiabatic_closure}
\end{align}
Equation~\eqref{eq:mass_conservation} conserves mass. Equation~\eqref{eq:mhd_momentum} is Newton's second law for the fluid. Equation~\eqref{eq:induction} evolves the magnetic field in an ideal conductor. Equation~\eqref{eq:divB} states that magnetic monopoles are absent. Equation~\eqref{eq:ampere_mhd} is Amp\`ere's law without the displacement-current term, appropriate when characteristic speeds are much smaller than the speed of light. Equation~\eqref{eq:adiabatic_closure} is one possible pressure closure, with $\gamma_{\mathrm{ad}}$ the adiabatic index.

The material derivative is
\begin{equation}
\frac{d}{dt}=\frac{\partial}{\partial t}+\mathbf u\cdot\nabla.
\end{equation}
It measures the rate of change experienced by a fluid element moving with velocity $\mathbf u$.

\subsection{Assumptions used for the baseline equilibrium}
The baseline Module~1 equilibrium adopts the following conceptual model unless a future implementation explicitly states otherwise:
\begin{enumerate}[leftmargin=1.6em]
\item \textbf{Time independence:} $\partial/\partial t=0$.
\item \textbf{No equilibrium bulk flow:} $\mathbf u=\mathbf 0$.
\item \textbf{Scalar pressure:} the pressure is represented by one scalar field $p$, rather than separate parallel and perpendicular pressures.
\item \textbf{Ideal magnetostatic force balance:} inertia, viscosity, and resistive electric-field effects are neglected in the equilibrium equation.
\item \textbf{Nested surfaces for the primary model:} the confined region is represented by smooth, nonintersecting magnetic surfaces. Islands and stochastic regions are discussed as limitations rather than included in the baseline representation.
\end{enumerate}
With these assumptions, Equation~\eqref{eq:mhd_momentum} reduces to the static equilibrium equation developed next.

\section{Magnetostatic equilibrium equations}
\label{sec:equilibrium_equations}

\subsection{The three governing relations}
The mathematical anchor for Module~1 is
\begin{align}
\nabla p&=\mathbf J\times\mathbf B,
\label{eq:force_balance}\\
\nabla\times\mathbf B&=\mu_0\mathbf J,
\label{eq:ampere}\\
\nabla\cdot\mathbf B&=0.
\label{eq:solenoidal}
\end{align}
Equation~\eqref{eq:force_balance} balances the pressure-gradient force against the magnetic Lorentz force density. Equation~\eqref{eq:ampere} relates current density to magnetic-field curl. Equation~\eqref{eq:solenoidal} requires the magnetic field to be divergence free.

The units provide an immediate check. Pressure gradient has units
\begin{equation}
[\nabla p]=\mathrm{Pa\,m^{-1}}=\mathrm{N\,m^{-3}}.
\end{equation}
The Lorentz force density has units
\begin{equation}
[\mathbf J\times\mathbf B]
=\left(\mathrm{A\,m^{-2}}\right)\left(\mathrm{T}\right)
=\mathrm{N\,m^{-3}},
\end{equation}
so the two sides of Equation~\eqref{eq:force_balance} are dimensionally consistent.

\subsection{Pressure is constant along a magnetic field line}
Take the dot product of Equation~\eqref{eq:force_balance} with $\mathbf B$:
\begin{equation}
\mathbf B\cdot\nabla p
=\mathbf B\cdot(\mathbf J\times\mathbf B)=0.
\label{eq:Bgradp_zero}
\end{equation}
A scalar field changes along a field line at a rate proportional to $\mathbf B\cdot\nabla p$. Therefore Equation~\eqref{eq:Bgradp_zero} states that pressure is constant along each magnetic field line.

If a field line densely covers a smooth magnetic surface, continuity implies that pressure is constant over the entire surface. We may then write
\begin{equation}
p=p(\psi),
\label{eq:p_flux_function}
\end{equation}
where $\psi$ labels magnetic surfaces. A quantity that depends only on the surface label is called a \textbf{flux function}.

Taking the dot product of Equation~\eqref{eq:force_balance} with $\mathbf J$ similarly gives
\begin{equation}
\mathbf J\cdot\nabla p=0.
\end{equation}
Thus both magnetic field lines and current lines lie tangent to constant-pressure surfaces in this ideal nested-surface equilibrium.

\subsection{Perpendicular and parallel current}
Decompose current density into components perpendicular and parallel to $\mathbf B$:
\begin{equation}
\mathbf J=\mathbf J_\perp+\mathbf J_\parallel,
\qquad
\mathbf J_\parallel=\frac{\mathbf J\cdot\mathbf B}{B^2}\mathbf B.
\end{equation}
Cross Equation~\eqref{eq:force_balance} with $\mathbf B$:
\begin{equation}
\mathbf B\times\nabla p
=\mathbf B\times(\mathbf J\times\mathbf B).
\end{equation}
Using the vector identity
\begin{equation}
\mathbf A\times(\mathbf C\times\mathbf D)
=\mathbf C(\mathbf A\cdot\mathbf D)-\mathbf D(\mathbf A\cdot\mathbf C),
\end{equation}
we obtain
\begin{equation}
\mathbf B\times(\mathbf J\times\mathbf B)
=B^2\mathbf J-(\mathbf B\cdot\mathbf J)\mathbf B
=B^2\mathbf J_\perp.
\end{equation}
Therefore
\begin{equation}
\boxed{\mathbf J_\perp=\frac{\mathbf B\times\nabla p}{B^2}.}
\label{eq:Jperp}
\end{equation}
The perpendicular current is directly fixed by pressure gradients and magnetic geometry. The parallel current is not determined by Equation~\eqref{eq:force_balance} alone; it is constrained together with geometry by Amp\`ere's law and charge conservation.

\subsection{Magnetic pressure and magnetic tension}
Substitute Equation~\eqref{eq:ampere} into the magnetic force density:
\begin{equation}
\mathbf J\times\mathbf B
=\frac{1}{\mu_0}(\nabla\times\mathbf B)\times\mathbf B.
\end{equation}
Using $\nabla\cdot\mathbf B=0$, this can be written as
\begin{equation}
\mathbf J\times\mathbf B
=-\nabla\left(\frac{B^2}{2\mu_0}\right)
+\frac{1}{\mu_0}(\mathbf B\cdot\nabla)\mathbf B.
\label{eq:magnetic_force_decomposition}
\end{equation}
The first term is the gradient of \textbf{magnetic pressure}, $B^2/(2\mu_0)$. The second term is \textbf{magnetic tension}, associated with bending field lines.

Writing $\mathbf B=B\mathbf b$ gives
\begin{equation}
(\mathbf B\cdot\nabla)\mathbf B
=B(\mathbf b\cdot\nabla B)\mathbf b+B^2\boldsymbol\kappa,
\end{equation}
where
\begin{equation}
\boldsymbol\kappa=\mathbf b\cdot\nabla\mathbf b
\label{eq:curvature}
\end{equation}
is the magnetic curvature vector. Perpendicular force balance can consequently be interpreted as a competition among plasma pressure, magnetic-pressure variation, and field-line tension.

\begin{physicsbox}
\textbf{Physical picture of equilibrium.} Plasma pressure tends to expand the plasma. Magnetic-pressure gradients and the tension of curved field lines provide restoring forces. A valid equilibrium is a three-dimensional arrangement in which these forces cancel at every point inside the plasma.
\end{physicsbox}

\subsection{Plasma beta}
A local pressure-to-magnetic-pressure ratio is
\begin{equation}
\beta(\mathbf x)=\frac{p(\mathbf x)}{B^2(\mathbf x)/(2\mu_0)}
=\frac{2\mu_0p(\mathbf x)}{B^2(\mathbf x)}.
\label{eq:local_beta}
\end{equation}
A volume-averaged beta may be defined by
\begin{equation}
\langle\beta\rangle_V
=\frac{2\mu_0\int_V p\,dV}{\int_V B^2\,dV}.
\label{eq:volume_beta}
\end{equation}
Different communities use slightly different averages, so a reported beta value must include its definition. Increasing beta strengthens the influence of plasma pressure on the magnetic geometry and can shift the magnetic axis, rotational transform, and stability properties.

\subsection{Equilibrium does not guarantee stability}
Equation~\eqref{eq:force_balance} states only that the net force vanishes for the chosen state. It does not state that a displaced plasma element experiences a restoring force. Stability requires perturbing the equilibrium and determining whether perturbations oscillate, decay, or grow. Module~4 begins that process for reduced shear-Alfv\'en waves. In symbols,
\begin{equation}
\text{equilibrium: }\mathbf F(\mathbf U_0)=0,
\qquad
\text{stability: behavior of }\delta\mathbf U\text{ near }\mathbf U_0.
\end{equation}
Here $\mathbf U_0$ denotes the equilibrium state and $\delta\mathbf U$ denotes a small perturbation.


% Module 1 section file. Included by ../module1_stellarator_equilibrium_geometry.tex.
\section{Coordinate systems for a toroidal plasma}
\label{sec:coordinates}

\subsection{Cartesian and cylindrical coordinates}
Cartesian coordinates $(x,y,z)$ use fixed orthonormal basis vectors. For a toroidal device, cylindrical coordinates $(R,\phi,Z)$ are more natural:
\begin{equation}
x=R\cos\phi,
\qquad
y=R\sin\phi,
\qquad z=Z.
\end{equation}
The coordinate $R\geq 0$ is distance from the vertical symmetry axis, $\phi$ is the geometric toroidal angle, and $Z$ is vertical position.

The cylindrical basis vectors $\mathbf e_R$, $\mathbf e_\phi$, and $\mathbf e_Z$ vary with $\phi$. A vector field may be written
\begin{equation}
\mathbf B=B_R\mathbf e_R+B_\phi\mathbf e_\phi+B_Z\mathbf e_Z.
\end{equation}
In an axisymmetric device, scalar quantities are independent of $\phi$. In a stellarator, this simplification generally does not hold.

\subsection{Major radius, minor radius, and aspect ratio}
The \textbf{magnetic axis} is the central closed field line around which nested magnetic surfaces are organized. A characteristic major radius $R_0$ measures the distance from the device's central axis to the magnetic axis. A characteristic minor radius $a$ measures the radial size of the plasma cross-section. The aspect ratio is
\begin{equation}
A=\frac{R_0}{a}.
\end{equation}
Large aspect ratio means the torus is relatively thin compared with its major radius. Reduced analytic models often exploit $a/R_0\ll1$, but realistic stellarator equilibria need not be accurately represented by this limit.

\subsection{Poloidal and toroidal angles}
A geometric poloidal angle may be introduced around a chosen cross-sectional center, while the geometric toroidal angle is $\phi$. Magnetic calculations often use modified angles $(\theta,\zeta)$ chosen to simplify field-line equations rather than to match geometric angles exactly. It is therefore essential to state whether an angle is geometric, computational, or a straight-field-line magnetic angle.

Both angular coordinates are periodic:
\begin{equation}
\theta\equiv\theta+2\pi,
\qquad
\zeta\equiv\zeta+2\pi.
\end{equation}
If the device has $N_{\mathrm{fp}}$ identical field periods, physical geometry repeats after
\begin{equation}
\Delta\zeta=\frac{2\pi}{N_{\mathrm{fp}}}.
\end{equation}

\section{Magnetic field lines and flux surfaces}
\label{sec:flux_surfaces}

\subsection{Definition of a magnetic field line}
A magnetic field line is a curve $\mathbf x(\ell)$ whose tangent is parallel to $\mathbf B$. One convenient parameter is arc length $\ell$, for which
\begin{equation}
\frac{d\mathbf x}{d\ell}=\mathbf b(\mathbf x)=\frac{\mathbf B(\mathbf x)}{B(\mathbf x)}.
\label{eq:field_line_arclength}
\end{equation}
An alternative parameter $\tau$ may be used:
\begin{equation}
\frac{d\mathbf x}{d\tau}=\mathbf B(\mathbf x).
\label{eq:field_line_B}
\end{equation}
Changing the parameter changes how quickly the curve is traversed, not the geometric path.

Field lines are mathematical integral curves. They are not material strings, although the magnetic-tension analogy is often useful. In ideal MHD, flux freezing gives them a strong dynamical connection to the plasma motion.

\subsection{Definition of a magnetic surface}
A magnetic surface is a surface labeled by a scalar $\psi(\mathbf x)$ such that
\begin{equation}
\mathbf B\cdot\nabla\psi=0.
\label{eq:B_tangent_surface}
\end{equation}
The gradient $\nabla\psi$ is normal to a constant-$\psi$ surface. Equation~\eqref{eq:B_tangent_surface} therefore says that $\mathbf B$ is tangent to the surface at every point. A field line that begins on the surface remains on it.

For nested surfaces, the plasma region resembles a collection of nonintersecting toroidal shells:
\begin{equation}
0\leq s\leq1,
\end{equation}
where $s=0$ labels the magnetic axis and $s=1$ labels the chosen plasma boundary, commonly the last closed flux surface.

\subsection{Toroidal and poloidal magnetic flux}
Magnetic flux through an oriented surface $S$ is
\begin{equation}
\Psi=\int_S\mathbf B\cdot d\mathbf S,
\end{equation}
where $d\mathbf S=\mathbf n\,dS$ and $\mathbf n$ is the surface normal.

For a toroidal magnetic surface, two fluxes are useful:
\begin{itemize}[leftmargin=1.6em]
\item The \textbf{toroidal flux} $\Psi_t$ passes through a poloidal cross-section of the torus.
\item The \textbf{poloidal flux} $\Psi_p$ passes through a toroidal ribbon extending around the long direction.
\end{itemize}
Sign conventions vary. This project will always store the convention explicitly rather than inferring it from a variable name.

A common normalized surface coordinate is
\begin{equation}
s=\frac{\Psi_t}{\Psi_{t,a}},
\label{eq:s_normalized_flux}
\end{equation}
where $\Psi_{t,a}$ is the toroidal flux at the selected plasma boundary. Another common radial coordinate is
\begin{equation}
\rho=\sqrt{s}.
\label{eq:rho_sqrt_s}
\end{equation}
For a simple circular plasma with nearly uniform toroidal field, toroidal flux scales approximately with the square of minor radius, so $\rho$ behaves approximately like a normalized geometric radius. In shaped three-dimensional geometry, $\rho$ remains a flux label rather than a literal Euclidean distance.

\subsection{Closed, islanded, and stochastic field-line regions}
The nested-surface picture is powerful but not universal.
\begin{itemize}[leftmargin=1.6em]
\item On an \textbf{irrational surface}, a field line generally does not close after a finite number of turns and can densely cover the surface.
\item On a \textbf{rational surface}, the field-line winding ratio is rational; perturbations can generate magnetic islands.
\item In a \textbf{stochastic region}, neighboring field lines separate chaotically and a single smooth global surface label may not exist.
\end{itemize}
A nested-surface equilibrium code such as standard VMEC represents the first category and idealized rational surfaces through smooth surfaces, but does not directly resolve magnetic islands or stochastic volumes. Other equilibrium formulations are needed when those structures are essential.


% Module 1 section file. Included by ../module1_stellarator_equilibrium_geometry.tex.
\section{Rotational transform, safety factor, and magnetic shear}
\label{sec:iota}

\subsection{Rotational transform}
Follow a field line and count how its poloidal angle changes as the toroidal angle increases. In straight-field-line coordinates, the rotational transform is
\begin{equation}
\iota(s)=\frac{d\theta}{d\zeta}
\label{eq:iota_dtheta_dzeta}
\end{equation}
along a field line, with the derivative understood as an average winding rate on the surface. Equivalently,
\begin{equation}
\iota(s)=\lim_{N\rightarrow\infty}\frac{\Delta\theta}{\Delta\zeta}
\end{equation}
for a trajectory followed through many toroidal turns.

If $\iota=0.4$, the field line advances on average $0.4$ poloidal turns per toroidal turn. After five toroidal turns it advances two poloidal turns. If $\iota=2/5$ exactly, the field line is rational in the idealized straight-field-line representation and closes after a finite number of turns.

\subsection{Safety factor}
Tokamak literature often uses the safety factor
\begin{equation}
q(s)=\frac{1}{\iota(s)}.
\label{eq:q_inverse_iota}
\end{equation}
With this convention, $q$ counts toroidal turns per poloidal turn. Some sign conventions include a minus sign depending on angle orientation and magnetic-field direction. This project uses Equation~\eqref{eq:q_inverse_iota} unless a data source documents a different convention.

\subsection{Rational surfaces}
A rational surface satisfies
\begin{equation}
\iota(s_r)=\frac{n}{m},
\label{eq:rational_iota}
\end{equation}
where $m$ and $n$ are nonzero integers with no common factor and $s_r$ is the rational surface location. A field line then closes after $m$ toroidal-angle periods and $n$ poloidal-angle periods under the stated convention. The same integers appear in Fourier perturbations, resonances, and Alfv\'en-continuum conditions.

\subsection{Magnetic shear}
Magnetic shear describes radial change in field-line pitch. Several definitions are used. A common dimensionless form in terms of $\iota$ is
\begin{equation}
\hat s_\iota(s)=\frac{s}{\iota(s)}\frac{d\iota}{ds}.
\label{eq:shear_iota}
\end{equation}
In terms of safety factor,
\begin{equation}
\hat s_q(s)=\frac{s}{q(s)}\frac{dq}{ds}=-\hat s_\iota(s).
\label{eq:shear_q}
\end{equation}
The sign difference follows from $q=1/\iota$. One must therefore report the definition, not only the phrase ``magnetic shear.''

Shear matters because nearby field lines accumulate different angular phases. It affects magnetic-island overlap, wave continua, resonant surfaces, and radial mode structure.

\subsection{Connection to the Module 4 parallel wave number}
Consider a perturbation with phase
\begin{equation}
S=m\theta-n\zeta-\omega t,
\label{eq:mode_phase}
\end{equation}
where $m$ is the poloidal mode number, $n$ is the toroidal mode number, and $\omega$ is angular frequency. Along a field line,
\begin{equation}
\frac{dS}{d\zeta}=m\frac{d\theta}{d\zeta}-n=m\iota-n.
\end{equation}
The parallel wave number is the phase gradient along the field:
\begin{equation}
k_\parallel=\mathbf b\cdot\nabla S.
\label{eq:kparallel_general}
\end{equation}
In a simple large-aspect-ratio approximation, $\mathbf b\cdot\nabla\zeta\approx1/R_0$, so
\begin{equation}
k_{\parallel,mn}(s)\approx\frac{m\iota(s)-n}{R_0}.
\label{eq:kparallel_simple}
\end{equation}
Module~4 may use the algebraically equivalent opposite sign $(n-m\iota)/R_0$ because its eigenfrequency depends on $k_\parallel^2$. The sign becomes important when propagation direction and resonant particle motion are retained.

\begin{physicsbox}
\textbf{Why Module 1 controls Module 4.} The equilibrium supplies $B(s,\theta,\zeta)$, mass density, metric information, and $\iota(s)$. These determine the local Alfv\'en speed and the parallel wavelength of each harmonic. Consequently, changing the equilibrium changes the Alfv\'en continuum and the locations at which energetic particles can resonate with waves.
\end{physicsbox}


% Module 1 section file. Included by ../module1_stellarator_equilibrium_geometry.tex.
\section{Magnetic flux coordinates}
\label{sec:magnetic_coordinates}

\subsection{Coordinate mapping}
Let
\begin{equation}
(u^1,u^2,u^3)=(s,\theta,\zeta)
\end{equation}
be curvilinear coordinates. The physical position is a mapping
\begin{equation}
\mathbf x=\mathbf x(s,\theta,\zeta).
\label{eq:coordinate_mapping}
\end{equation}
The covariant basis vectors are
\begin{equation}
\mathbf e_i=\frac{\partial\mathbf x}{\partial u^i},
\end{equation}
and the reciprocal, or contravariant, basis vectors are
\begin{equation}
\mathbf e^i=\nabla u^i.
\end{equation}
They satisfy
\begin{equation}
\mathbf e_i\cdot\mathbf e^j=\delta_i^{\ j},
\end{equation}
where $\delta_i^{\ j}$ is the Kronecker delta, equal to one when $i=j$ and zero otherwise.

\subsection{Metric tensors}
The covariant metric tensor is
\begin{equation}
g_{ij}=\mathbf e_i\cdot\mathbf e_j,
\label{eq:covariant_metric}
\end{equation}
and the contravariant metric tensor is
\begin{equation}
g^{ij}=\mathbf e^i\cdot\mathbf e^j.
\label{eq:contravariant_metric}
\end{equation}
The matrices $[g_{ij}]$ and $[g^{ij}]$ are inverses. They convert coordinate derivatives and components into physical lengths and dot products.

For example, the squared differential distance is
\begin{equation}
d\ell^2=g_{ij}\,du^i du^j,
\label{eq:line_element}
\end{equation}
where repeated indices are summed from 1 to 3.

\subsection{Jacobian and volume element}
The coordinate Jacobian is
\begin{equation}
\sqrt{g}=\mathbf e_s\cdot(\mathbf e_\theta\times\mathbf e_\zeta).
\label{eq:jacobian}
\end{equation}
The physical volume element is
\begin{equation}
dV=\sqrt{g}\,ds\,d\theta\,d\zeta.
\label{eq:volume_element}
\end{equation}
The sign of $\sqrt{g}$ depends on coordinate orientation. A production interface should either enforce a positive orientation or store the sign convention explicitly.

\subsection{Contravariant representation of a tangent, divergence-free field}
A particularly useful straight-field-line representation is
\begin{equation}
\boxed{
\mathbf B
=\nabla\psi\times\nabla\theta
+\iota(\psi)\,\nabla\zeta\times\nabla\psi.
}
\label{eq:B_contravariant}
\end{equation}
Here $\psi$ is a toroidal-flux-like surface coordinate. Equation~\eqref{eq:B_contravariant} automatically gives
\begin{equation}
\mathbf B\cdot\nabla\psi=0
\end{equation}
and, provided $\iota$ depends only on $\psi$, satisfies $\nabla\cdot\mathbf B=0$.

Using the Jacobian identity
\begin{equation}
(\nabla\psi\times\nabla\theta)\cdot\nabla\zeta=\frac{1}{\sqrt{g}},
\end{equation}
we find
\begin{align}
\mathbf B\cdot\nabla\zeta&=\frac{1}{\sqrt{g}},\\
\mathbf B\cdot\nabla\theta&=\frac{\iota}{\sqrt{g}}.
\end{align}
Therefore
\begin{equation}
\frac{d\theta}{d\zeta}
=\frac{\mathbf B\cdot\nabla\theta}{\mathbf B\cdot\nabla\zeta}
=\iota,
\end{equation}
which explains the phrase \textbf{straight-field-line coordinates}: in the $(\theta,\zeta)$ plane, field lines are straight lines of slope $\iota$.

\subsection{Covariant representation and Boozer coordinates}
Any vector can also be expanded in gradients:
\begin{equation}
\mathbf B=B_s\nabla s+B_\theta\nabla\theta+B_\zeta\nabla\zeta.
\end{equation}
In Boozer coordinates the covariant representation can be written
\begin{equation}
\mathbf B
=K(s,\theta,\zeta)\nabla s
+I(s)\nabla\theta
+G(s)\nabla\zeta.
\label{eq:B_boozer_covariant}
\end{equation}
The functions $I(s)$ and $G(s)$ are flux functions related to net toroidal and poloidal currents, while $K$ may vary over a surface. Exact multiplicative and sign conventions depend on how flux and angles are normalized. This is why a geometry file must store metadata rather than only arrays.

Boozer coordinates are useful because the magnetic-field-line trajectories are straight and the covariant angular components are flux functions. Guiding-center orbit equations and quasisymmetry are especially transparent in this representation. Detailed treatments of flux coordinates and the original Boozer construction are given in Refs.~\cite{dHaeseleer1991,boozer1981}.

\subsection{Hamada coordinates}
Hamada coordinates are another straight-field-line system. They are constructed so that both magnetic field lines and current lines have simple contravariant representations and the Jacobian is a flux function under suitable conditions. Boozer and Hamada coordinates serve different analytical and numerical purposes; neither is simply a geometric angle system.

\subsection{Coordinate singularity at the magnetic axis}
At $s=0$, all poloidal angles refer to the same physical magnetic axis. The coordinate system is therefore singular in the same way that polar angle is undefined at the origin. This is a coordinate singularity, not necessarily a physical singularity. Numerical representations must impose regularity conditions so that physical fields remain finite and single-valued near the axis.


% Module 1 section file. Included by ../module1_stellarator_equilibrium_geometry.tex.
\section{Flux-surface averages and radial geometry}
\label{sec:surface_averages}

\subsection{Surface average}
For a scalar field $A(s,\theta,\zeta)$, a common flux-surface average is
\begin{equation}
\langle A\rangle_s
=\frac{\displaystyle\int_0^{2\pi}\int_0^{2\pi}A\sqrt{g}\,d\theta\,d\zeta}
{\displaystyle\int_0^{2\pi}\int_0^{2\pi}\sqrt{g}\,d\theta\,d\zeta}.
\label{eq:flux_surface_average}
\end{equation}
The Jacobian appears because equal increments of $(\theta,\zeta)$ do not generally correspond to equal physical surface or volume weighting.

\subsection{Enclosed volume and differential volume}
The volume inside surface $s$ is
\begin{equation}
V(s)=\int_0^s ds'\int_0^{2\pi}d\theta\int_0^{2\pi}d\zeta\,\sqrt{g}(s',\theta,\zeta).
\label{eq:enclosed_volume}
\end{equation}
Its derivative is
\begin{equation}
V'(s)=\frac{dV}{ds}
=\int_0^{2\pi}\int_0^{2\pi}\sqrt{g}\,d\theta\,d\zeta.
\label{eq:Vprime}
\end{equation}
The function $V'(s)$ appears in one-dimensional transport equations. For example, a conservative radial balance often contains
\begin{equation}
\frac{1}{V'}\frac{\partial}{\partial s}\left(V'\Gamma_s\right),
\end{equation}
where $\Gamma_s$ is a radial flux measured per appropriate surface area and coordinate normalization.

\subsection{Effective minor radius}
A coordinate-independent size measure may be defined from volume:
\begin{equation}
a_{\mathrm{eff}}(s)=\sqrt{\frac{V(s)}{2\pi^2R_{\mathrm{ref}}}},
\end{equation}
where $R_{\mathrm{ref}}$ is a stated reference major radius. This definition is motivated by the volume $2\pi^2R_0a^2$ of a circular torus. Other effective-radius definitions exist, so the formula must accompany the value.

\section{Three-dimensional stellarator surface geometry}
\label{sec:fourier_geometry}

\subsection{Why use Fourier series?}
The angles $\theta$ and $\zeta$ are periodic. Any sufficiently smooth periodic surface shape can therefore be expanded in Fourier modes. A common representation in cylindrical coordinates is
\begin{align}
R(s,\theta,\zeta)
&=\sum_{m,n}R_{mn}(s)
\cos\left(m\theta-nN_{\mathrm{fp}}\zeta\right),
\label{eq:R_fourier}\\
Z(s,\theta,\zeta)
&=\sum_{m,n}Z_{mn}(s)
\sin\left(m\theta-nN_{\mathrm{fp}}\zeta\right).
\label{eq:Z_fourier}
\end{align}
Here $m$ is a poloidal integer, $n$ is a toroidal Fourier integer per field period, and $N_{\mathrm{fp}}$ is the number of field periods. The coefficients $R_{mn}(s)$ and $Z_{mn}(s)$ determine the shape of every flux surface.

The physical position is then
\begin{equation}
\mathbf x(s,\theta,\zeta)
=R\cos\zeta\,\mathbf e_x
+R\sin\zeta\,\mathbf e_y
+Z\,\mathbf e_z,
\label{eq:x_from_RZ}
\end{equation}
when $\zeta$ is chosen as the geometric toroidal angle. If $\zeta$ is a magnetic toroidal angle, an additional mapping to the geometric angle is required.

\subsection{Meaning of selected Fourier modes}
The coefficient $R_{00}$ largely sets a major-radius offset. The $m=1,n=0$ terms describe a basic circular or elliptical cross-sectional displacement. Higher poloidal mode numbers describe triangularity, squareness, and finer shaping. Nonzero toroidal mode numbers introduce variation from one toroidal location to another and are essential to a stellarator.

A Fourier coefficient should not be interpreted in isolation without knowing the angular convention. The same physical surface can have different coefficients after a change of poloidal angle.

\subsection{Stellarator symmetry}
A common discrete symmetry is stellarator symmetry. Under an appropriate choice of angular origins, the configuration is invariant under
\begin{equation}
(\theta,\zeta)\rightarrow(-\theta,-\zeta)
\end{equation}
combined with a reflection in physical space. This symmetry motivates cosine terms for $R$ and sine terms for $Z$ in Equations~\eqref{eq:R_fourier}--\eqref{eq:Z_fourier}. Configurations without stellarator symmetry require both sine and cosine coefficients in both coordinates.

\subsection{Magnetic-field-strength spectrum}
The magnetic-field magnitude can also be expanded:
\begin{equation}
B(s,\theta,\zeta)
=\sum_{m,n}B_{mn}(s)
\cos\left(m\theta-nN_{\mathrm{fp}}\zeta+\delta_{mn}(s)\right),
\label{eq:B_spectrum_general}
\end{equation}
where $\delta_{mn}$ is a phase. With suitable symmetry and angle conventions, the phases can be absorbed into sine/cosine parity choices.

The surface shape spectrum and magnetic-strength spectrum are related through Maxwell's equations and equilibrium, but they are not identical. Particle trapping depends strongly on variation of $B$ along a field line, so the $B_{mn}$ spectrum is a critical Module~3 input.

\subsection{Spectral resolution and aliasing}
A numerical representation retains finite ranges
\begin{equation}
0\leq m\leq m_{\max},
\qquad
-n_{\max}\leq n\leq n_{\max}.
\end{equation}
Increasing $m_{\max}$ and $n_{\max}$ resolves finer structure but increases computational cost and can amplify sensitivity to noisy input. Nonlinear products generate sums of mode numbers, so evaluation on an insufficient angular grid can cause aliasing: unresolved high-frequency content appears incorrectly as lower-frequency content. Production workflows should document spectral truncation and quadrature resolution.


% Module 1 section file. Included by ../module1_stellarator_equilibrium_geometry.tex.
\section{Vacuum magnetic field and external coils}
\label{sec:coils}

\subsection{Biot--Savart law for a filamentary coil}
A current-carrying wire contributes a magnetic field
\begin{equation}
\mathbf B(\mathbf x)
=\frac{\mu_0 I_c}{4\pi}
\oint_C
\frac{d\boldsymbol\ell'\times(\mathbf x-\mathbf x')}
{|\mathbf x-\mathbf x'|^3},
\label{eq:biot_savart}
\end{equation}
where $I_c$ is coil current, $C$ is the coil centerline, $\mathbf x'$ is a source point on the coil, and $d\boldsymbol\ell'$ is an oriented line element. Fields from multiple coils superpose linearly in vacuum.

Real coils have finite cross-section, support structures, current-density variation, and manufacturing tolerances. A filament model is often adequate for conceptual studies but must be refined for engineering analysis near the conductor.

\subsection{Vacuum equations}
In a current-free vacuum region,
\begin{equation}
\nabla\times\mathbf B=0,
\qquad
\nabla\cdot\mathbf B=0.
\end{equation}
Locally one may introduce a magnetic scalar potential $\Phi_m$ such that
\begin{equation}
\mathbf B=-\nabla\Phi_m,
\qquad
\nabla^2\Phi_m=0,
\end{equation}
provided the region is simply connected or global topological fluxes are treated separately. Alternatively, a vector potential $\mathbf A$ may be used:
\begin{equation}
\mathbf B=\nabla\times\mathbf A,
\end{equation}
which automatically enforces $\nabla\cdot\mathbf B=0$.

\subsection{Coil field versus plasma response}
The total field inside an operating plasma is
\begin{equation}
\mathbf B_{\mathrm{total}}
=\mathbf B_{\mathrm{coils}}+\mathbf B_{\mathrm{plasma}},
\label{eq:total_field}
\end{equation}
where $\mathbf B_{\mathrm{plasma}}$ is generated by equilibrium plasma currents. At low beta and low net plasma current, the coil field may dominate. At finite beta, diamagnetic and parallel currents modify the geometry. A vacuum field-line calculation is therefore not automatically a finite-pressure MHD equilibrium.

\section{Boundary conditions and equilibrium classes}
\label{sec:boundary_conditions}

\subsection{Fixed-boundary equilibrium}
In a fixed-boundary calculation, the plasma boundary shape is prescribed:
\begin{equation}
\mathbf x_a(\theta,\zeta)=\mathbf x(s=1,\theta,\zeta).
\end{equation}
The boundary is required to be a magnetic surface,
\begin{equation}
\mathbf B\cdot\mathbf n=0
\qquad\text{on }s=1,
\label{eq:fixed_boundary_Bn}
\end{equation}
where $\mathbf n$ is the outward unit normal. The solver determines the interior surfaces and field consistent with the pressure and current profiles.

Fixed-boundary calculations are useful when a target plasma shape is already known or when equilibrium is embedded inside an optimization loop that treats boundary shape as the design variable.

\subsection{Free-boundary equilibrium}
In a free-boundary calculation, external coils and currents are prescribed, and the plasma boundary is solved as part of the coupled plasma-vacuum problem. The interface must satisfy normal-field and stress-balance conditions. In the idealized case of a tangential magnetic field at the boundary, total pressure balance is
\begin{equation}
\left[
 p+\frac{B^2}{2\mu_0}
\right]_{\mathrm{inside}}
=
\left[
\frac{B^2}{2\mu_0}
\right]_{\mathrm{outside}},
\label{eq:total_pressure_boundary}
\end{equation}
where the outside is vacuum and square brackets here indicate evaluation on the two sides of the interface. Surface currents may permit a jump in tangential magnetic field, so a complete formulation must state whether such currents are allowed.

\subsection{Magnetic axis regularity}
At the magnetic axis, physical quantities must remain finite and single-valued. Fourier components must scale with radius in a manner consistent with regularity. For example, non-axisymmetric poloidal harmonics should vanish with appropriate powers of the minor-radius coordinate as $s\rightarrow0$. Numerical solvers enforce this through basis choices, parity conditions, or explicit constraints.

\subsection{Profile boundary conditions}
Pressure typically decreases from a central value to a small edge value:
\begin{equation}
p(0)=p_0,
\qquad
p(1)\approx0.
\end{equation}
The exact edge condition is a model choice. A current profile may be specified through toroidal current, parallel current, or derivatives of flux functions. Profiles must be differentiable enough for the chosen numerical method; sharp edge gradients require adequate resolution.

\section{Pressure, density, temperature, and current profiles}
\label{sec:profiles}

\subsection{Pressure profile}
A simple analytic pressure profile is
\begin{equation}
p(s)=p_0(1-s)^{\alpha_p},
\label{eq:pressure_profile_example}
\end{equation}
where $p_0$ is central pressure and $\alpha_p>0$ controls peaking. Its derivative is
\begin{equation}
\frac{dp}{ds}=-\alpha_p p_0(1-s)^{\alpha_p-1}.
\end{equation}
The pressure gradient, not only the pressure value, drives perpendicular current through Equation~\eqref{eq:Jperp}.

For multiple species with isotropic temperatures,
\begin{equation}
p(s)=\sum_s n_s(s)k_{\mathrm B}T_s(s),
\label{eq:multispecies_pressure}
\end{equation}
where $k_{\mathrm B}$ is Boltzmann's constant. If temperature is specified in electronvolts, the conversion to joules must be handled consistently.

\subsection{Mass density profile}
The MHD mass density is
\begin{equation}
\rho_m(s)=\sum_s m_s n_s(s).
\label{eq:mass_density_species}
\end{equation}
Electron mass is often negligible in the mass sum but not in charge balance or kinetic physics. Module~4 uses $\rho_m$ to compute the Alfv\'en speed
\begin{equation}
v_A=\frac{B}{\sqrt{\mu_0\rho_m}}.
\label{eq:alfven_speed_module1}
\end{equation}
Thus a geometry product that omits density is insufficient for a fully physical continuum calculation, even though density may be supplied by Module~7 rather than solved by Module~1.

\subsection{Current-profile measures}
In a stellarator, several current measures can be relevant:
\begin{itemize}[leftmargin=1.6em]
\item net toroidal plasma current;
\item bootstrap current driven by collisional transport in pressure gradients;
\item externally driven current from neutral beams or radio-frequency waves;
\item Pfirsch--Schluter currents required to maintain current continuity in three-dimensional geometry;
\item diamagnetic current associated with pressure gradients.
\end{itemize}
These categories are not independent arbitrary additions; they are different physical contributions to the total current density that must remain consistent with force balance and Ampere's law.

\subsection{Profile consistency}
An equilibrium solver cannot accept an arbitrary collection of mutually incompatible functions. Pressure, current, toroidal flux, and rotational transform are linked by the equilibrium equations. Depending on the solver formulation, one specifies some profiles and solves for the others. The input contract must state which quantities are prescribed and which are outputs.

\begin{notebox}
\textbf{Common modeling error.} Supplying a desired pressure profile, a desired rotational-transform profile, and a desired current profile does not guarantee that all three can coexist with a prescribed boundary. An equilibrium code enforces compatibility; it is not merely an interpolation tool.
\end{notebox}


% Module 1 section file. Included by ../module1_stellarator_equilibrium_geometry.tex.
\section{Axisymmetry versus fully three-dimensional equilibrium}
\label{sec:axisymmetry_vs_3d}

\subsection{Axisymmetric reduction}
An axisymmetric equilibrium is invariant under toroidal rotation:
\begin{equation}
\frac{\partial}{\partial\phi}=0.
\end{equation}
This continuous symmetry produces a conserved canonical toroidal momentum for particle motion and permits the magnetic field to be represented using a two-dimensional poloidal-flux function. The equilibrium equations reduce to the Grad--Shafranov equation, a nonlinear elliptic partial differential equation in the $(R,Z)$ plane.

The exact Grad--Shafranov form depends on flux normalization. A common form is
\begin{equation}
R\frac{\partial}{\partial R}
\left(\frac{1}{R}\frac{\partial\psi_p}{\partial R}\right)
+\frac{\partial^2\psi_p}{\partial Z^2}
=-\mu_0R^2\frac{dp}{d\psi_p}
-F\frac{dF}{d\psi_p},
\label{eq:grad_shafranov}
\end{equation}
where $\psi_p(R,Z)$ is a poloidal-flux function and $F(\psi_p)=RB_\phi$. This equation is included to show what axisymmetry buys mathematically: a scalar two-dimensional equilibrium equation.

\subsection{Why a stellarator is harder}
A stellarator generally has only discrete toroidal periodicity, not continuous axisymmetry. Quantities depend on all three coordinates:
\begin{equation}
A=A(s,\theta,\zeta).
\end{equation}
There is no single two-dimensional Grad--Shafranov equation that describes a general three-dimensional stellarator equilibrium. Instead, numerical methods solve coupled equations for the shapes of many nested surfaces and for field/current quantities over each surface.

Three-dimensionality affects more than computational cost. It changes the confinement physics summarized in the review by Helander~\cite{helander2014}:
\begin{itemize}[leftmargin=1.6em]
\item magnetic-well structure and trapped-particle orbits;
\item bootstrap and Pfirsch--Schluter currents;
\item neoclassical transport;
\item energetic-particle confinement;
\item coupling among poloidal and toroidal wave harmonics;
\item coil complexity and sensitivity to errors.
\end{itemize}

\section{Quasisymmetry and related optimization ideas}
\label{sec:quasisymmetry}

\subsection{Why symmetry matters for particle confinement}
In an exactly axisymmetric magnetic field, the particle Lagrangian is independent of the toroidal angle. Noether's theorem then gives conservation of canonical toroidal momentum. This extra invariant strongly restricts radial particle excursions.

A stellarator lacks exact axisymmetry in physical space, but it can be designed so that the magnetic-field magnitude $B$ has an approximate continuous symmetry in Boozer coordinates. This property is called \textbf{quasisymmetry}. The vector field remains three-dimensional, while the magnitude sampled by guiding centers has a simpler dependence.

\subsection{General quasisymmetric form}
A quasisymmetric magnetic-field strength has the form
\begin{equation}
B=B\left(s,M\theta-N\zeta\right),
\label{eq:quasisymmetry_form}
\end{equation}
where $M$ and $N$ are fixed integers that define the symmetry direction. The most common cases are:
\begin{itemize}[leftmargin=1.6em]
\item \textbf{quasi-axisymmetry:} $N=0$, so $B$ is approximately independent of the Boozer toroidal angle;
\item \textbf{quasi-helical symmetry:} both $M$ and $N$ are nonzero, so $B$ approximately follows a helical angle.
\end{itemize}
The word ``quasi'' is essential: exact global quasisymmetry is severely constrained, and practical configurations minimize symmetry-breaking Fourier components rather than eliminating all of them. Modern high-precision constructions are discussed by Landreman and Paul~\cite{landremanPaul2022}.

\subsection{Quasi-isodynamic configurations}
Another stellarator optimization strategy is quasi-isodynamicity, associated with favorable contours of the second adiabatic invariant and reduced radial drift of trapped particles. It is not the same condition as quasisymmetry. Both strategies aim to improve confinement but organize the magnetic geometry differently.

\subsection{Why Module 1 should not reduce geometry to one score}
A single quasisymmetry-error metric cannot fully characterize an equilibrium. A useful geometry database should retain the field spectrum, rotational transform, magnetic shear, curvature, magnetic wells, effective ripple measures, and coil/plasma constraints. Different downstream modules are sensitive to different features.

\section{Geometric quantities needed for particle orbits}
\label{sec:orbit_geometry}

\subsection{Field magnitude and unit vector}
Guiding-center motion depends on
\begin{equation}
B=|\mathbf B|,
\qquad
\mathbf b=\frac{\mathbf B}{B}.
\end{equation}
These must be available as continuous functions or sufficiently accurate interpolants in physical space or magnetic coordinates.

\subsection{Magnetic-field gradient}
The gradient
\begin{equation}
\nabla B
\end{equation}
drives the grad-$B$ drift. A noisy numerical derivative can cause artificial orbit diffusion even when $B$ itself is accurate. Geometry export should therefore include either analytically differentiable spectral representations or derivative fields computed with controlled accuracy.

\subsection{Field-line curvature}
The curvature
\begin{equation}
\boldsymbol\kappa=\mathbf b\cdot\nabla\mathbf b
\end{equation}
drives curvature drift and contributes to MHD interchange physics. Its dimensions are inverse length. The curvature vector is perpendicular to $\mathbf b$ because differentiating $\mathbf b\cdot\mathbf b=1$ along the field gives
\begin{equation}
\mathbf b\cdot\boldsymbol\kappa=0.
\end{equation}

\subsection{Mirror ratio and magnetic wells}
On a chosen field line or surface, define
\begin{equation}
B_{\min}(s)=\min_{\theta,\zeta}B(s,\theta,\zeta),
\qquad
B_{\max}(s)=\max_{\theta,\zeta}B(s,\theta,\zeta).
\end{equation}
The mirror ratio is
\begin{equation}
\mathcal R_m(s)=\frac{B_{\max}(s)}{B_{\min}(s)}.
\end{equation}
Particles with sufficiently large perpendicular velocity can be reflected from high-$B$ regions and become trapped. The detailed arrangement of local minima and maxima along a field line matters more than the mirror ratio alone.

\subsection{Electric potential}
Module 1 primarily supplies magnetic equilibrium, but guiding-center orbits may also require an electrostatic potential
\begin{equation}
\Phi=\Phi(s,\theta,\zeta),
\qquad
\mathbf E=-\nabla\Phi
\end{equation}
for a static electrostatic field. A lowest-order model often assumes $\Phi=\Phi(s)$, giving a radial electric field. Because this potential is not determined by ideal magnetostatic force balance alone, its source and coupling location must be documented separately.

\section{Geometric quantities needed for waves and stability}
\label{sec:wave_geometry}

\subsection{Parallel derivative}
The derivative along the magnetic field is
\begin{equation}
\nabla_\parallel A=\mathbf b\cdot\nabla A
\end{equation}
for a scalar $A$. In magnetic coordinates, the contravariant field representation provides
\begin{equation}
\mathbf B\cdot\nabla A
=\frac{1}{\sqrt g}
\left(\frac{\partial A}{\partial\zeta}
+\iota\frac{\partial A}{\partial\theta}\right)
\label{eq:Bdotgrad_straight}
\end{equation}
under the convention of Equation~\eqref{eq:B_contravariant}. This operator is central to shear-Alfven propagation.

\subsection{Local Alfven frequency}
For a Fourier harmonic proportional to $\exp[i(m\theta-n\zeta-\omega t)]$, a reduced local shear-Alfven frequency is
\begin{equation}
\omega_A(s,\theta,\zeta)
\sim |k_\parallel|v_A,
\qquad
v_A=\frac{B}{\sqrt{\mu_0\rho_m}}.
\end{equation}
In a three-dimensional equilibrium, $B$, the metric, and sometimes density vary over a surface, so a full continuum calculation is more involved than the simple radial expression used in the reduced Module 4 benchmark.

\subsection{Curvature and pressure-gradient coupling}
MHD stability operators contain combinations of pressure gradient, current, field-line curvature, and metric coefficients. Bad curvature means that pressure forces tend to displace plasma farther in the direction of the original displacement; good curvature tends to provide restoration. In a three-dimensional torus, curvature varies strongly along a field line, so stability depends on weighted averages and mode structure rather than a single local sign.

\subsection{Geodesic curvature}
Curvature can be decomposed relative to a flux surface into normal and tangential components. Geodesic curvature lies within the surface and is perpendicular to the field line. It participates in coupling between shear-Alfven and compressional or acoustic responses. A production equilibrium interface should provide enough metric and derivative information to calculate this decomposition consistently.


% Module 1 section file. Included by ../module1_stellarator_equilibrium_geometry.tex.
\section{A reduced analytic equilibrium for intuition}
\label{sec:analytic_model}

\subsection{Circular nested surfaces}
Before using a three-dimensional equilibrium, it is useful to define a manufactured geometry with circular concentric surfaces:
\begin{align}
R(r,\theta)&=R_0+r\cos\theta,\\
Z(r,\theta)&=r\sin\theta,
\end{align}
where $0\leq r\leq a$. Let
\begin{equation}
s=\left(\frac{r}{a}\right)^2,
\qquad
r=a\sqrt{s}.
\end{equation}
This model is axisymmetric and therefore not a stellarator, but it provides an interpretable verification limit for coordinate mappings, Jacobians, volume, and interpolation.

The enclosed volume is approximately
\begin{equation}
V(r)=2\pi^2R_0r^2,
\end{equation}
and hence
\begin{equation}
V(s)=2\pi^2R_0a^2s,
\qquad
V'(s)=2\pi^2R_0a^2.
\end{equation}
These relations give exact or near-exact checks for a large-aspect-ratio circular implementation.

\subsection{Prescribed rotational transform}
A simple transform profile is
\begin{equation}
\iota(s)=\iota_0+\iota_1s,
\label{eq:iota_linear}
\end{equation}
where $\iota_0$ is the on-axis transform and $\iota_1$ is the total edge-minus-axis change. The shear measure becomes
\begin{equation}
\hat s_\iota(s)
=\frac{s\iota_1}{\iota_0+\iota_1s}.
\end{equation}
This profile is not derived from force balance; it is a manufactured geometry input suitable for checking downstream equations.

\subsection{Simple helical field-line tracing}
In straight-field-line coordinates,
\begin{equation}
\frac{d\theta}{d\zeta}=\iota(s).
\end{equation}
For constant $s$ and constant $\iota$,
\begin{equation}
\theta(\zeta)=\theta_0+\iota\zeta,
\label{eq:field_line_solution}
\end{equation}
where $\theta_0$ is the initial poloidal angle. This exact solution is a basic field-line tracer test.

\subsection{What the analytic model omits}
The circular manufactured model omits three-dimensional shaping, finite-beta equilibrium shifts, realistic current distributions, magnetic wells, coil ripple, islands, and stochasticity. It is a verification case, not a predictive stellarator equilibrium.


% Module 1 section file. Included by ../module1_stellarator_equilibrium_geometry.tex.
\section{Numerical solution of three-dimensional equilibrium}
\label{sec:equilibrium_solvers}

\subsection{Unknowns and input choices}
A nested-surface equilibrium calculation typically represents the position of each surface by spectral coefficients such as $R_{mn}(s)$ and $Z_{mn}(s)$. Depending on formulation, it also solves for angular mappings, magnetic-field components, current-related flux functions, and the rotational-transform profile.

Typical fixed-boundary inputs include:
\begin{itemize}[leftmargin=1.6em]
\item plasma boundary Fourier coefficients;
\item toroidal magnetic flux;
\item pressure profile;
\item either current information or rotational-transform constraints;
\item radial and angular resolution;
\item symmetry and sign conventions.
\end{itemize}
Typical free-boundary inputs additionally include coil geometry, coil currents, conducting structures, and vacuum-region settings.

\subsection{Variational principle}
Many ideal-MHD equilibrium methods use an energy principle. A representative total energy is
\begin{equation}
W=\int_V\left[
\frac{B^2}{2\mu_0}
+\frac{p}{\gamma_{\mathrm{ad}}-1}
\right]dV,
\label{eq:equilibrium_energy}
\end{equation}
subject to constraints such as conservation of magnetic flux and an equation of state. A stationary point of this constrained functional satisfies ideal-MHD force balance.

The phrase \textbf{variational} means that the solver changes the surface representation to reduce a residual or energy gradient derived from $W$. It does not mean that every stationary point is automatically stable; the second variation is relevant to stability.

\subsection{VMEC concept}
The Variational Moments Equilibrium Code, abbreviated VMEC, represents nested flux surfaces spectrally and seeks a three-dimensional ideal-MHD equilibrium by minimizing force residuals associated with a variational formulation. The radial direction is discretized, and angular dependence is represented using Fourier series.

Important modeling assumptions include:
\begin{enumerate}[leftmargin=1.6em]
\item nested toroidal flux surfaces are imposed;
\item pressure is a flux function;
\item magnetic islands and stochastic regions are not directly represented;
\item convergence depends on radial and spectral resolution, initial guess, profiles, and boundary shape.
\end{enumerate}
VMEC output is often transformed into Boozer coordinates using a separate post-processing step because the coordinates used internally by the equilibrium solver are not automatically the coordinates best suited for orbit and neoclassical calculations. The foundational variational moment method is described by Hirshman and Whitson~\cite{hirshmanWhitson1983}.

\subsection{Alternative equilibrium formulations}
Other formulations include force-balance Newton solvers, near-axis expansions, free-boundary integral methods, and multi-region relaxed MHD. Multi-region models can represent interfaces between relaxed plasma volumes and can accommodate magnetic islands or stochastic regions more naturally than a globally nested-surface model~\cite{hudson2012}. The appropriate formulation depends on the physics question.

\subsection{Discretization}
A typical numerical representation uses:
\begin{itemize}[leftmargin=1.6em]
\item a radial grid $s_j$, with index $j=0,\ldots,N_s-1$ for $N_s$ stored radial points;
\item Fourier modes $(m,n)$ up to selected truncations;
\item numerical quadrature in the angular directions;
\item finite differences, finite elements, splines, or spectral bases in radius.
\end{itemize}
The unknown coefficient vector may be written abstractly as
\begin{equation}
\mathbf a=
\{R_{mn}(s_j),Z_{mn}(s_j),\ldots\}.
\end{equation}
The discrete nonlinear equilibrium equations are
\begin{equation}
\mathbf F(\mathbf a)=\mathbf 0,
\label{eq:nonlinear_equilibrium_discrete}
\end{equation}
where $\mathbf F$ is a vector of force-balance and constraint residuals.

\subsection{Nonlinear iteration}
A Newton method updates
\begin{equation}
\mathbf a^{(k+1)}
=\mathbf a^{(k)}+\delta\mathbf a^{(k)},
\end{equation}
where
\begin{equation}
\mathbf J_F(\mathbf a^{(k)})\delta\mathbf a^{(k)}
=-\mathbf F(\mathbf a^{(k)}),
\end{equation}
and $\mathbf J_F$ is the Jacobian matrix of the residual. Variational descent methods instead move along a preconditioned negative energy gradient. Practical codes use damping, preconditioning, continuation in pressure or boundary deformation, and problem-specific regularization.

\subsection{Continuation}
A difficult finite-beta equilibrium may be reached gradually. Introduce a continuation parameter $0\leq\lambda\leq1$ and solve a sequence
\begin{equation}
p_\lambda(s)=\lambda p_{\mathrm{target}}(s).
\end{equation}
The converged solution at one $\lambda$ becomes the initial guess for the next. Continuation can also be applied to boundary shaping, plasma current, or coil perturbations.

\subsection{Convergence criteria}
A solver should not declare convergence from iteration count alone. Useful criteria include:
\begin{align}
\epsilon_F&=\frac{\|\nabla p-\mathbf J\times\mathbf B\|}{F_{\mathrm{ref}}},\\
\epsilon_B&=\frac{\|\nabla\cdot\mathbf B\|}{B_{\mathrm{ref}}/L_{\mathrm{ref}}},\\
\epsilon_a&=\frac{\|\mathbf a^{(k+1)}-\mathbf a^{(k)}\|}{\|\mathbf a^{(k)}\|+a_{\mathrm{floor}}}.
\end{align}
Here $F_{\mathrm{ref}}$, $B_{\mathrm{ref}}$, $L_{\mathrm{ref}}$, and $a_{\mathrm{floor}}$ are explicitly chosen scaling quantities. Norm type, quadrature weighting, and excluded singular points must be documented.

\subsection{Force residuals in curvilinear coordinates}
A small scalar residual does not guarantee that every physical component is small. The force residual
\begin{equation}
\mathbf f=\mathbf J\times\mathbf B-\nabla p
\end{equation}
should be examined in radial, poloidal, and toroidal directions, or decomposed into components parallel and perpendicular to $\mathbf B$. Surface-localized residuals can be hidden by a global average.

\subsection{Resolution studies}
Convergence must be checked with respect to:
\begin{itemize}[leftmargin=1.6em]
\item radial grid size $N_s$;
\item maximum poloidal mode $m_{\max}$;
\item maximum toroidal mode $n_{\max}$;
\item angular quadrature grid;
\item nonlinear tolerance.
\end{itemize}
A geometry quantity $Q$ should approach a stable value as resolution increases. A useful relative difference is
\begin{equation}
\delta_Q^{(h)}
=\frac{|Q_h-Q_{h/2}|}{|Q_{h/2}|+Q_{\mathrm{floor}}},
\end{equation}
where $h$ denotes an effective discretization scale and $Q_{\mathrm{floor}}$ prevents division by a negligible number.

\subsection{Failure modes}
Common equilibrium-solver failures include:
\begin{itemize}[leftmargin=1.6em]
\item self-intersecting or poorly parameterized surfaces;
\item loss of nested surfaces in the physical configuration;
\item insufficient spectral resolution;
\item excessive pressure gradient or incompatible profiles;
\item coordinate singularities and negative Jacobian values;
\item convergence to an unintended branch;
\item sensitivity to initial guess;
\item free-boundary mismatch between plasma and vacuum fields.
\end{itemize}
A failed solve must produce explicit flags and diagnostics rather than silently exporting geometry.


% Module 1 section file. Included by ../module1_stellarator_equilibrium_geometry.tex.
\section{Coordinate transformation and interpolation}
\label{sec:transform_interpolation}

\subsection{Why transformations are unavoidable}
Different calculations prefer different representations:
\begin{itemize}[leftmargin=1.6em]
\item coil and wall geometry are naturally expressed in Cartesian or cylindrical coordinates;
\item equilibrium solvers use internal flux coordinates and spectral coefficients;
\item guiding-center solvers often prefer Boozer coordinates or direct Cartesian fields;
\item wave solvers require metric coefficients and parallel derivatives;
\item transport solvers frequently use a one-dimensional normalized flux coordinate.
\end{itemize}
Module 1 must therefore define transformations, not only raw field values.

\subsection{Forward and inverse mappings}
The forward mapping
\begin{equation}
(s,\theta,\zeta)\mapsto\mathbf x
\end{equation}
is evaluated from the surface representation. The inverse mapping
\begin{equation}
\mathbf x\mapsto(s,\theta,\zeta)
\end{equation}
is generally nonlinear and may require root finding. The inverse can be nonunique outside the nested-surface region or near self-intersections, so its domain must be restricted.

A Newton inverse solves
\begin{equation}
\mathbf r(s,\theta,\zeta)
=\mathbf x(s,\theta,\zeta)-\mathbf x_{\mathrm{target}}
=\mathbf 0.
\end{equation}
The update uses the coordinate-basis matrix
\begin{equation}
\begin{bmatrix}
\partial_s\mathbf x & \partial_\theta\mathbf x & \partial_\zeta\mathbf x
\end{bmatrix}.
\end{equation}
Near a coordinate singularity or a poorly shaped surface, this matrix can be ill-conditioned.

\subsection{Interpolation requirements}
A downstream solver may evaluate geometry at points not present on the equilibrium grid. Interpolation should preserve:
\begin{itemize}[leftmargin=1.6em]
\item periodicity in angular coordinates;
\item continuity across field-period boundaries;
\item regularity at the magnetic axis;
\item stated order of accuracy;
\item consistency between a field and its derivatives.
\end{itemize}
Interpolating $B$, $\nabla B$, and $\boldsymbol\kappa$ independently can violate derivative identities. A differentiable spectral representation is preferable when practical.

\subsection{Conservative remapping of flux functions}
When a profile such as density or pressure is transferred from one radial grid to another, pointwise interpolation does not necessarily conserve integrated particle number or energy. A conservative remap should preserve an integral such as
\begin{equation}
\mathcal N
=\int_0^1 n(s)V'(s)\,ds.
\label{eq:particle_integral}
\end{equation}
For cell averages, this is accomplished by matching integrals over overlapping radial cells rather than simply matching values at cell centers.


% Module 1 section file. Included by ../module1_stellarator_equilibrium_geometry.tex.
\section{Module 1 output contract}
\label{sec:outputs}

\subsection{Minimum geometry product}
A useful equilibrium data product must identify the equilibrium uniquely and contain enough information for independent checks. At minimum, it should include:
\begin{enumerate}[leftmargin=1.6em]
\item coordinate definitions and angle orientation;
\item unit system and normalization constants;
\item field-period number $N_{\mathrm{fp}}$;
\item magnetic-axis representation;
\item surface mapping $\mathbf x(s,\theta,\zeta)$ or equivalent Fourier coefficients;
\item magnetic field $\mathbf B$ and magnitude $B$;
\item rotational transform $\iota(s)$;
\item pressure $p(s)$ and, when used, mass density $\rho_m(s)$;
\item Jacobian $\sqrt g$ and metric tensors or enough information to reconstruct them;
\item enclosed volume $V(s)$ and $V'(s)$;
\item derivatives required for $\nabla B$, curvature, and parallel operators;
\item boundary and last-closed-surface definition;
\item convergence diagnostics and failure flags.
\end{enumerate}

\subsection{Recommended metadata}
Recommended metadata include:
\begin{itemize}[leftmargin=1.6em]
\item solver name and version;
\item source input-file checksum;
\item fixed-boundary or free-boundary designation;
\item pressure and current-profile formulas;
\item radial and spectral resolution;
\item nonlinear tolerance and achieved residuals;
\item sign convention for toroidal flux, poloidal flux, $\iota$, and angles;
\item whether $\zeta$ is geometric or magnetic toroidal angle;
\item coordinate transformation method;
\item date and provenance of coil geometry;
\item known warnings, extrapolation regions, and invalid domains.
\end{itemize}

\subsection{Suggested array shapes}
For a sampled grid with $N_s$ radial points, $N_\theta$ poloidal points, and $N_\zeta$ toroidal points, representative arrays are
\begin{align}
\texttt{position}&:[N_s,N_\theta,N_\zeta,3],\\
\texttt{B\_vector}&:[N_s,N_\theta,N_\zeta,3],\\
\texttt{B\_magnitude}&:[N_s,N_\theta,N_\zeta],\\
\texttt{jacobian}&:[N_s,N_\theta,N_\zeta],\\
\texttt{iota}&:[N_s],\\
\texttt{pressure}&:[N_s],\\
\texttt{Vprime}&:[N_s].
\end{align}
The final software schema may use spectral coefficients instead of dense arrays. Whatever representation is chosen, axis order and memory layout must be documented.

\subsection{Downstream consumers}
The main consumers are:
\begin{longtable}{>{\RaggedRight\arraybackslash}p{0.16\textwidth} >{\RaggedRight\arraybackslash}p{0.30\textwidth} >{\RaggedRight\arraybackslash}p{0.44\textwidth}}
\toprule
\textbf{Module} & \textbf{Required Module 1 quantities} & \textbf{Physical use}\\
\midrule
\endfirsthead
\toprule
\textbf{Module} & \textbf{Required Module 1 quantities} & \textbf{Physical use}\\
\midrule
\endhead
Module 2 & $s$, volume, profiles, field strength & Places energetic-particle sources and distributions on the equilibrium.\\
Module 3 & $\mathbf B$, $B$, $\nabla B$, $\boldsymbol\kappa$, coordinates, wall mapping & Advances guiding centers and identifies confined and lost trajectories.\\
Module 4 & $\iota$, $B$, $\rho_m$, metric, $\mathbf b\cdot\nabla$ & Computes shear-Alfven continua, harmonic coupling, and eigenmodes.\\
Module 5 & mode geometry plus orbit frequencies derived from Module 1 & Evaluates wave-particle resonances and energetic-particle drive.\\
Module 6 & $V'$, radial coordinate, mode geometry & Evolves reduced fast-ion transport conservatively.\\
Module 7 & $V'$, surface areas, density and field data & Evolves heating and thermal profiles.\\
Module 8 & all typed state products and provenance & Coordinates multiphysics exchanges and detects stale geometry.\\
Module 9 & parameterized equilibrium inputs and physical residuals & Trains and validates equilibrium or downstream surrogate models.\\
Module 10 & viewing geometry and field mapping & Generates synthetic diagnostic signals.\\
Module 11 & named physical dependencies and interventions & Learns candidate coupling graphs without treating arrays as anonymous features.\\
\bottomrule
\caption{Principal equilibrium dependencies across the project.}
\label{tab:module1_consumers}
\end{longtable}

\subsection{Versioning and immutable state products}
An equilibrium should be treated as an immutable state product with a unique identifier. If pressure, coil current, boundary shape, or plasma current changes enough to require a new solve, the resulting geometry receives a new identifier. Downstream outputs should record the exact equilibrium identifier used. This prevents a mode spectrum or orbit-loss map from being paired accidentally with a different geometry.


% Module 1 section file. Included by ../module1_stellarator_equilibrium_geometry.tex.
\section{Worked conceptual examples}
\label{sec:worked_examples}

\subsection{Example 1: pressure as a flux function}
Suppose a nested magnetic surface is labeled by $s$. Start from force balance,
\begin{equation}
\nabla p=\mathbf J\times\mathbf B.
\end{equation}
Dot with $\mathbf B$:
\begin{equation}
\mathbf B\cdot\nabla p=0.
\end{equation}
Parameterize a field line by $\tau$ using $d\mathbf x/d\tau=\mathbf B$. The pressure derivative along the line is
\begin{equation}
\frac{dp}{d\tau}
=\frac{d\mathbf x}{d\tau}\cdot\nabla p
=\mathbf B\cdot\nabla p
=0.
\end{equation}
Therefore $p$ is constant along that field line. If the line densely covers the surface and $p$ is continuous, all points on the surface have the same pressure, so $p=p(s)$.

\subsection{Example 2: field-line closure}
Let $\iota=2/5$ and use the straight-line solution
\begin{equation}
\theta(\zeta)=\theta_0+\frac{2}{5}\zeta.
\end{equation}
After five toroidal turns, $\Delta\zeta=10\pi$. The poloidal advance is
\begin{equation}
\Delta\theta=\frac{2}{5}(10\pi)=4\pi,
\end{equation}
which is two poloidal turns. The angular coordinates return to their starting values modulo $2\pi$, so the ideal field line closes.

\subsection{Example 3: a rational Alfven resonance surface}
Take a mode with $(m,n)=(2,1)$ and the transform profile
\begin{equation}
\iota(s)=0.35+0.30s.
\end{equation}
The simple parallel wave number vanishes when
\begin{equation}
2\iota(s)-1=0.
\end{equation}
Thus
\begin{equation}
2(0.35+0.30s)-1=0
\quad\Rightarrow\quad
s=0.5.
\end{equation}
The surface $s=0.5$ is the rational surface $\iota=1/2$ for this harmonic. In a more complete shear-Alfven model, coupling, compressibility, and finite-frequency effects determine the actual continuum behavior near this location.

\subsection{Example 4: beta estimate}
Let $p=2.0\times10^5\ \mathrm{Pa}$ and $B=2.5\ \mathrm T$. The magnetic pressure is
\begin{equation}
\frac{B^2}{2\mu_0}
=\frac{(2.5)^2}{2(4\pi\times10^{-7})}
\approx2.49\times10^6\ \mathrm{Pa}.
\end{equation}
The local beta is
\begin{equation}
\beta\approx\frac{2.0\times10^5}{2.49\times10^6}
\approx0.080,
\end{equation}
or approximately $8.0\%$. This does not mean magnetic forces exceed pressure forces everywhere by exactly a factor of $12.5$; equilibrium depends on gradients and curvature, not only local energy-density ratios.

\subsection{Example 5: field-period Fourier variation}
Suppose $N_{\mathrm{fp}}=4$ and a surface contains the term
\begin{equation}
R_{1,1}(s)\cos(\theta-4\zeta).
\end{equation}
Increasing $\zeta$ by one field period, $\Delta\zeta=2\pi/4$, changes the phase by
\begin{equation}
-4\Delta\zeta=-2\pi.
\end{equation}
The cosine is unchanged, confirming fourfold toroidal periodicity.


% Module 1 section file. Included by ../module1_stellarator_equilibrium_geometry.tex.
\section{Verification and validation}
\label{sec:verification}

\subsection{Definitions}
\textbf{Verification} asks whether the equations were solved correctly. \textbf{Validation} asks whether those equations and inputs adequately represent the physical system for the intended use. A solver can be verified but invalid for a given experiment, or physically appropriate but implemented incorrectly.

\subsection{Code verification tests}
A Module 1 implementation should include at least the following tests.

\subsubsection{Dimensional and convention tests}
Every exported quantity must carry units or an explicit normalization. Tests should check angle ranges, coordinate orientation, flux signs, and the relation $q=1/\iota$ under the project convention.

\subsubsection{Divergence-free test}
Evaluate
\begin{equation}
\epsilon_{\nabla\cdot B}
=\frac{L_{\mathrm{ref}}\|\nabla\cdot\mathbf B\|_2}{\|\mathbf B\|_2}.
\end{equation}
The norm should decrease with increased numerical resolution until limited by solver tolerance or interpolation error.

\subsubsection{Surface-tangency test}
Evaluate
\begin{equation}
\epsilon_{B\cdot\nabla s}
=\frac{\|\mathbf B\cdot\nabla s\|_2}
{\|\mathbf B\|_\infty\|\nabla s\|_2+\epsilon_{\mathrm{floor}}}.
\end{equation}
A nested-surface equilibrium should make this small throughout the valid domain.

\subsubsection{Force-balance test}
Evaluate
\begin{equation}
\epsilon_{\mathrm{force}}
=\frac{\|\mathbf J\times\mathbf B-\nabla p\|_2}
{\|\nabla p\|_2+\|\mathbf J\times\mathbf B\|_2+\epsilon_{\mathrm{floor}}}.
\end{equation}
The test should also inspect maximum and surface-resolved residuals.

\subsubsection{Field-line transform test}
Trace field lines independently and estimate
\begin{equation}
\iota_{\mathrm{trace}}(s)
=\frac{\Delta\theta}{\Delta\zeta}
\end{equation}
over many turns. Compare with the exported $\iota(s)$. The discrepancy should decrease with tracing tolerance and interpolation resolution.

\subsubsection{Flux conservation test}
Compute magnetic flux through equivalent cross-sections of the same flux tube. Because $\nabla\cdot\mathbf B=0$, the flux should be independent of cross-section up to numerical error.

\subsubsection{Coordinate round-trip test}
For points inside the valid domain, apply
\begin{equation}
(s,\theta,\zeta)\rightarrow\mathbf x
\rightarrow(\tilde s,\tilde\theta,\tilde\zeta).
\end{equation}
Compare $s$ directly and compare angles modulo $2\pi$. Failures near the axis should be treated with coordinate-aware tolerances.

\subsubsection{Manufactured circular-torus test}
Recover the analytic mappings, volume, $V'$, and field-line solution of Section~\ref{sec:analytic_model}. This separates geometry-interface errors from the complexity of a full equilibrium code.

\subsection{Solution verification}
Even a verified code requires case-specific convergence studies. For each equilibrium, report changes in quantities such as
\begin{itemize}[leftmargin=1.6em]
\item magnetic-axis location;
\item $\iota(s)$;
\item volume-averaged beta;
\item minimum Jacobian;
\item selected $B_{mn}$ coefficients;
\item $B_{\min}$ and $B_{\max}$;
\item force residual;
\item derived Module 4 continuum frequencies.
\end{itemize}
The last item is a cross-module sensitivity test: a geometry may appear visually converged while small transform or metric errors still shift resonant frequencies.

\subsection{Validation sources}
Possible validation sources include:
\begin{itemize}[leftmargin=1.6em]
\item comparison with an independently implemented equilibrium code;
\item comparison with published benchmark equilibria;
\item magnetic-diagnostic reconstruction from an experiment;
\item pressure-profile and plasma-boundary measurements;
\item vacuum-field mapping before plasma operation;
\item cross-checks with coil-current perturbation experiments.
\end{itemize}
Validation must be tied to an intended prediction. Agreement in boundary shape does not automatically validate energetic-particle orbit confinement.

\subsection{Uncertainty sources}
Important uncertainties include:
\begin{itemize}[leftmargin=1.6em]
\item coil-position and coil-current errors;
\item plasma-pressure and current-profile uncertainty;
\item diagnostic reconstruction error;
\item finite spectral and radial resolution;
\item coordinate-transformation error;
\item model-form error from scalar-pressure ideal MHD;
\item omitted islands, stochasticity, flow, or anisotropy.
\end{itemize}
Uncertainty should be propagated to quantities actually used by downstream modules, such as $\iota$, orbit-loss fraction, or eigenfrequency.

\section{Equilibrium surrogates and physics constraints}
\label{sec:surrogates}

\subsection{Why build a surrogate?}
A high-fidelity equilibrium solve inside a large parameter scan, control loop, uncertainty analysis, or digital twin can be expensive. A surrogate seeks a rapid approximation
\begin{equation}
\mathcal G:
\left(\text{boundary, coils, profiles, controls}\right)
\mapsto
\left(\mathbf B,\iota,\text{geometry, diagnostics}\right).
\end{equation}
Possible models include neural operators, reduced-basis methods, Gaussian processes, and physics-informed neural networks.

\subsection{Physical constraints}
A useful equilibrium surrogate should be checked against
\begin{align}
\nabla\cdot\mathbf B&=0,\\
\mathbf B\cdot\nabla s&=0,\\
\mathbf J-\frac{1}{\mu_0}\nabla\times\mathbf B&=0,\\
\mathbf J\times\mathbf B-\nabla p&=0,
\end{align}
in the domain where these equations apply. It should also respect periodicity, stellarator symmetry when imposed, boundary conditions, and positive Jacobian orientation.

\subsection{Representation choices}
Predicting Cartesian components of $\mathbf B$ independently can violate $\nabla\cdot\mathbf B=0$. Better choices may include:
\begin{itemize}[leftmargin=1.6em]
\item predicting a vector potential $\mathbf A$ and setting $\mathbf B=\nabla\times\mathbf A$;
\item predicting flux-surface coordinates and using a divergence-free contravariant representation;
\item predicting spectral coefficients constrained by symmetry;
\item learning corrections to a reduced conventional equilibrium rather than the full field.
\end{itemize}
Each representation moves some physical constraints from the loss function into the architecture.

\subsection{Validation requirements}
Low average data error is insufficient. A surrogate should be tested for:
\begin{enumerate}[leftmargin=1.6em]
\item in-distribution interpolation;
\item out-of-distribution boundary and profile changes;
\item force and divergence residuals;
\item topology preservation and positive Jacobian;
\item accuracy of derived quantities such as $\iota$, $k_\parallel$, curvature, and orbit invariants;
\item uncertainty calibration;
\item graceful failure or rejection outside the training domain.
\end{enumerate}

\subsection{Do not replace the reference solver}
The conventional equilibrium solver remains the reference for generating training data, checking residuals, and identifying surrogate failure. The surrogate accelerates selected workflows; it does not eliminate the need for conventional physics verification.

\section{Limitations of the baseline Module 1 model}
\label{sec:limitations}

The textbook description and planned interface are broader than the current executable repository capability. The baseline does not yet implement a validated three-dimensional equilibrium solver. Furthermore, the primary physical model described here has the following limitations unless explicitly extended:
\begin{enumerate}[leftmargin=1.6em]
\item scalar, isotropic pressure rather than an anisotropic pressure tensor;
\item static equilibrium without strong bulk flow;
\item ideal MHD rather than resistive, two-fluid, or kinetic equilibrium;
\item globally nested flux surfaces;
\item no self-consistent magnetic islands or stochastic regions;
\item no scrape-off-layer or open-field-line equilibrium;
\item no direct engineering stress, thermal, or structural coil calculation;
\item no guarantee of MHD stability or good particle confinement;
\item no automatic electrostatic-potential solution;
\item no uncertainty-calibrated experimental reconstruction.
\end{enumerate}
These boundaries are not defects when stated clearly. They define the domain in which Module 1 outputs may be trusted.

\section{Conceptual map for interview-level explanation}
\label{sec:interview_map}

A concise but technically correct explanation is:

\begin{physicsbox}
A stellarator equilibrium is a three-dimensional arrangement of magnetic field, plasma current, and pressure that satisfies $\nabla p=\mathbf J\times\mathbf B$, $\nabla\times\mathbf B=\mu_0\mathbf J$, and $\nabla\cdot\mathbf B=0$. The field is organized into nested toroidal magnetic surfaces. Field lines wind poloidally and toroidally on each surface, with average pitch measured by the rotational transform $\iota$. Surface shape, field-strength variation, metric coefficients, curvature, and $\iota$ determine particle orbits, Alfven-wave propagation, resonances, and transport. A stellarator obtains this geometry primarily from three-dimensional external coils, so equilibrium computation and coordinate transformation are foundational inputs to every later kinetic-MHD module.
\end{physicsbox}

The following distinctions should remain explicit:
\begin{itemize}[leftmargin=1.6em]
\item equilibrium versus stability;
\item vacuum field versus finite-pressure plasma equilibrium;
\item geometric angles versus straight-field-line angles;
\item magnetic surface label versus physical radius;
\item rotational transform $\iota$ versus safety factor $q$;
\item accurate field values versus accurate derivatives;
\item a nested-surface representation versus a field containing islands or stochasticity;
\item a surrogate prediction versus a verified conventional solution.
\end{itemize}


% Module 1 section file. Included by ../module1_stellarator_equilibrium_geometry.tex.
\section{Chapter summary}
\label{sec:summary}

\begin{enumerate}[leftmargin=1.6em]
\item Magnetic confinement exploits rapid gyromotion and relatively free motion along magnetic field lines. Closing field lines toroidally avoids simple end loss.
\item A stellarator uses three-dimensional external magnetic geometry to produce field-line twist without relying on a large transformer-driven plasma current.
\item Static ideal-MHD equilibrium satisfies
\begin{equation*}
\nabla p=\mathbf J\times\mathbf B,
\qquad
\nabla\times\mathbf B=\mu_0\mathbf J,
\qquad
\nabla\cdot\mathbf B=0.
\end{equation*}
\item Force balance implies $\mathbf B\cdot\nabla p=0$, so pressure is a flux function when nested surfaces exist.
\item A magnetic surface satisfies $\mathbf B\cdot\nabla s=0$. The normalized toroidal-flux coordinate $s$ labels nested surfaces, while $\rho=\sqrt{s}$ is a convenient radius-like coordinate.
\item Rotational transform $\iota=d\theta/d\zeta$ measures average field-line pitch. The safety factor is $q=1/\iota$ under the project convention.
\item Radial variation of $\iota$ produces magnetic shear. Rational values $\iota=n/m$ identify surfaces important for islands and wave resonances.
\item Straight-field-line coordinates make field lines straight in angular coordinate space. Boozer coordinates are especially useful for particle-orbit and quasisymmetry analysis.
\item Three-dimensional surfaces and field strength are represented efficiently by Fourier series in poloidal and toroidal angles.
\item The vacuum coil field and plasma-current field both contribute to the operating equilibrium.
\item Fixed-boundary calculations prescribe the plasma boundary; free-boundary calculations solve the plasma-vacuum interface with coil fields.
\item A numerical equilibrium is a nonlinear constrained problem requiring radial and spectral convergence studies, force-balance checks, and coordinate-quality diagnostics.
\item Downstream modules require not only $\mathbf B$ but also $B$, derivatives, curvature, metrics, Jacobian, $\iota$, $V'$, profiles, conventions, provenance, and uncertainty information.
\item A physics-informed surrogate must preserve equilibrium constraints and derived-quantity accuracy; it cannot be validated by field-value error alone.
\end{enumerate}

\appendix

\section{Vector identities used in this chapter}
\label{app:vector_identities}

\subsection{Triple-product identities}
For vectors $\mathbf A$, $\mathbf C$, and $\mathbf D$,
\begin{align}
\mathbf A\cdot(\mathbf C\times\mathbf D)
&=\mathbf C\cdot(\mathbf D\times\mathbf A)
=\mathbf D\cdot(\mathbf A\times\mathbf C),\\
\mathbf A\times(\mathbf C\times\mathbf D)
&=\mathbf C(\mathbf A\cdot\mathbf D)
-\mathbf D(\mathbf A\cdot\mathbf C).
\end{align}

\subsection{Magnetic-force identity}
For a differentiable vector field $\mathbf B$,
\begin{equation}
(\nabla\times\mathbf B)\times\mathbf B
=(\mathbf B\cdot\nabla)\mathbf B
-\frac{1}{2}\nabla B^2
-\mathbf B(\nabla\cdot\mathbf B).
\end{equation}
When $\nabla\cdot\mathbf B=0$, this yields Equation~\eqref{eq:magnetic_force_decomposition}.

\subsection{Divergence of a curl}
For a sufficiently smooth vector potential $\mathbf A$,
\begin{equation}
\nabla\cdot(\nabla\times\mathbf A)=0.
\end{equation}
Thus representing $\mathbf B=\nabla\times\mathbf A$ enforces the solenoidal condition analytically.

\section{Derivation of the straight-field-line representation}
\label{app:straight_field_derivation}

Let $\psi$ label magnetic surfaces and let $\theta$ and $\zeta$ be periodic angles. A field tangent to surfaces can be represented using cross products containing $\nabla\psi$:
\begin{equation}
\mathbf B
=A(\psi,\theta,\zeta)\nabla\psi\times\nabla\theta
+C(\psi,\theta,\zeta)\nabla\zeta\times\nabla\psi.
\end{equation}
The two terms are tangent to constant-$\psi$ surfaces. A suitable angular-coordinate choice and flux normalization set $A=1$ and make $C=\iota(\psi)$, giving Equation~\eqref{eq:B_contravariant}. This coordinate construction is nontrivial globally; the equation should be understood as the defining representation of an appropriate straight-field-line system, not as a statement that arbitrary geometric angles already have this property.

\section{Example project interface checklist}
\label{app:interface_checklist}

Before a Module 1 state product is accepted by another module, verify:
\begin{enumerate}[leftmargin=1.6em]
\item The equilibrium identifier and source checksum are present.
\item SI units or normalization constants are present for every dimensional field.
\item The meanings of $s$, $\rho$, $\theta$, and $\zeta$ are documented.
\item Angle orientation, flux signs, and the $q$--$\iota$ convention are documented.
\item The valid domain and last closed flux surface are identified.
\item The field-period convention and Fourier convention are identified.
\item $\mathbf B$, $B$, position, Jacobian, metric, and $\iota$ are mutually consistent.
\item Derivatives are generated from the same representation as the primary field where possible.
\item The minimum Jacobian is positive under the stored orientation.
\item Force, divergence, tangency, and transform residuals pass specified tolerances.
\item Radial and spectral convergence evidence is attached.
\item Interpolation and extrapolation behavior is documented.
\item The geometry does not silently extrapolate outside the valid plasma or vacuum domain.
\item Downstream outputs retain the equilibrium identifier.
\end{enumerate}

\section{Review questions and exercises}
\label{app:exercises}

\begin{enumerate}[leftmargin=1.8em]
\item Explain why a purely toroidal magnetic field is not generally sufficient for good toroidal confinement. Distinguish particle streaming from cross-field drifts.
\item Starting from $\nabla p=\mathbf J\times\mathbf B$, prove that both $\mathbf B$ and $\mathbf J$ are tangent to constant-pressure surfaces.
\item Derive Equation~\eqref{eq:Jperp} for the perpendicular current.
\item Use the magnetic-force identity to interpret magnetic pressure and tension.
\item For $B=3\ \mathrm T$ and $p=1.5\times10^5\ \mathrm{Pa}$, calculate the local beta.
\item A field line has $\iota=0.42$. How many poloidal turns does it make during ten toroidal turns? Is the line guaranteed to close?
\item For $\iota(s)=0.3+0.4s$, locate the surfaces with $\iota=1/2$ and $\iota=2/3$, if they lie in $0\leq s\leq1$.
\item Show that Equation~\eqref{eq:B_contravariant} gives $d\theta/d\zeta=\iota$.
\item Explain why $\rho=\sqrt{s}$ should not automatically be interpreted as geometric distance from the magnetic axis.
\item Describe the difference between fixed-boundary and free-boundary equilibrium.
\item Explain why a vacuum coil field is not necessarily the same as a finite-beta plasma equilibrium.
\item What physics is excluded when globally nested surfaces are imposed?
\item Why can accurate interpolation of $B$ still produce inaccurate guiding-center orbits if $\nabla B$ is noisy?
\item Write the Fourier phase for a mode with poloidal number $m=3$, toroidal number $n=2$, and $N_{\mathrm{fp}}=4$ under the convention in Equation~\eqref{eq:R_fourier}.
\item Explain the distinction between stellarator symmetry and quasisymmetry.
\item Give three quantities from Module 1 that Module 4 needs and explain why.
\item Construct a convergence table for $\iota(0.5)$ as $m_{\max}$, $n_{\max}$, and $N_s$ are increased. What pattern would support numerical convergence?
\item Propose a divergence-free neural representation for $\mathbf B$ and identify one disadvantage.
\item Explain why a low mean-square error in $B$ does not guarantee a useful equilibrium surrogate.
\item State the most important assumptions that must accompany any claim based on this Module 1 model.
\end{enumerate}


% Module 1 section file. Included by ../module1_stellarator_equilibrium_geometry.tex.
\renewcommand{\refname}{References and further reading}
\phantomsection
\addcontentsline{toc}{section}{References and further reading}
\label{sec:references}
\begin{thebibliography}{99}

\bibitem{freidberg2014}
J. P. Freidberg,
\textit{Ideal MHD},
Cambridge University Press, Cambridge, 2014.

\bibitem{dHaeseleer1991}
W. D. D'haeseleer, W. N. G. Hitchon, J. D. Callen, and J. L. Shohet,
\textit{Flux Coordinates and Magnetic Field Structure: A Guide to a Fundamental Tool of Plasma Theory},
Springer, Berlin, 1991.

\bibitem{helander2014}
P. Helander,
``Theory of plasma confinement in non-axisymmetric magnetic fields,''
\textit{Reports on Progress in Physics}, vol. 77, 087001, 2014.

\bibitem{helanderSigmar2002}
P. Helander and D. J. Sigmar,
\textit{Collisional Transport in Magnetized Plasmas},
Cambridge University Press, Cambridge, 2002.

\bibitem{wesson2011}
J. Wesson,
\textit{Tokamaks}, 4th ed.,
Oxford University Press, Oxford, 2011.

\bibitem{boozer1981}
A. H. Boozer,
``Plasma equilibrium with rational magnetic surfaces,''
\textit{Physics of Fluids}, vol. 24, pp. 1999--2003, 1981.

\bibitem{hirshmanWhitson1983}
S. P. Hirshman and J. C. Whitson,
``Steepest-descent moment method for three-dimensional magnetohydrodynamic equilibria,''
\textit{Physics of Fluids}, vol. 26, pp. 3553--3568, 1983.

\bibitem{hudson2012}
S. R. Hudson, R. L. Dewar, G. Dennis, M. J. Hole, M. McGann, G. von Nessi, and S. Lazerson,
``Computation of multi-region relaxed magnetohydrodynamic equilibria,''
\textit{Physics of Plasmas}, vol. 19, 112502, 2012.

\bibitem{landremanPaul2022}
M. Landreman and E. Paul,
``Magnetic fields with precise quasisymmetry for plasma confinement,''
\textit{Physical Review Letters}, vol. 128, 035001, 2022.

\bibitem{hazeltineMeiss2003}
R. D. Hazeltine and J. D. Meiss,
\textit{Plasma Confinement},
Dover Publications, Mineola, New York, 2003.

\end{thebibliography}



\end{document}
